\documentclass[11pt, a4paper]{article}

% --- UNIVERSAL PREAMBLE BLOCK ---
\usepackage[a4paper, top=2.5cm, bottom=2.5cm, left=2cm, right=2cm]{geometry}
\usepackage{fontspec}

\usepackage[english, bidi=basic, provide=*]{babel}

\babelprovide[import, onchar=ids fonts]{english}

% Set default/Latin font to Sans Serif in the main (rm) slot
\babelfont{rm}{Noto Sans}

% --- STANDARD PACKAGES ---
\usepackage{amsmath}
\usepackage{amssymb}
\usepackage{amsthm}
\usepackage{booktabs}
\usepackage{longtable}
\usepackage{array}
\usepackage{hyperref}

\hypersetup{
    colorlinks=true,
    linkcolor=blue,
    urlcolor=blue,
    citecolor=blue,
    pdftitle={Mathematical Architecture and Interaction Mechanics of the Zeta Energy System (Hoard 2026i & Fifteen Lemmas Complete Edition)}
}

\newtheorem{proposition}{Proposition}[section]
\newtheorem{definition}{Definition}[section]
\newtheorem{theorem}{Theorem}[section]
\newtheorem{corollary}{Corollary}[section]
\newtheorem{lemma}{Lemma}[section]

\title{\textbf{Mathematical Architecture, Kinematic Calculations, and Interaction Mechanics of the Zeta Energy System}\\[0.4em]
\large A Complete Walk-Through of the Hoard (2026i) Geometric Engine of Time, the Fifteen Lemmas, the $720 \to 72 \to 10$ Torsional Markov Pump, Monoid Collapse, Golden Strassen-Clifford Toroids, and Rational Trigonometric Replacements}
\author{BioChain AI SHD-CCP Architecture Research \\ \small Technical Specification, Cross-Paper Synthesis \& Computational Audit}
\date{November 2026}

\begin{document}

\maketitle

\begin{abstract}
This paper provides a complete, equation-by-equation walk-through of the calculations implemented in the upgraded \texttt{zeta\_energy\_system.html} simulation environment and integrates the complete deductive foundation of four companion monographs: \textit{The Fifteen Lemmas of the Geometric Engine} (Hoard, October 16, 2026), \textit{Monoid Collapse and Mass-Node Traversal in Relational Space}, \textit{The Torsional Markov Memory-Offload Pump} (Hoard, June 21, 2026), and \textit{Quaternionic Delta Encoding on Golden Strassen-Clifford Toroids}. Every computational subsystem is derived in detail: (1) the 30-wheel prime-weighted Dirichlet/Zeta signal proposer and its $C_3$ Fourier decomposition into an apex scalar invariant $s$ and rotating plane $z$; (2) the projective Gaussian spinor clock $(1+i)^k$ in the binary octahedral group $2O$ and its square-root-free rational spread envelope $s_k = \sin^2(k\pi/4)$; (3) the joint 24-frame Chinese Remainder Theorem (CRT) schedule and its $720 \to 72 \to 10$ torsional gear nesting; (4) the embedding of the $\{4,3,5\}$ golden hyperbolic honeycomb cell, the corrected nine-phase tetrahedral wheel $\mu_9$ in $\mathbb{Q}(2\cos 40^\circ)$, and the $D_3$-symmetric trefoil knot; (5) the seven distinct mathematical interaction mechanisms that govern trajectory collisions, inter-slice bridges, conformal null-cone inner products $\langle X_A, X_B\rangle$ in $\mathbb{R}^{4,1}$, and 64-bit Einstein Tile telemetry; and (6) the complete integration of all Fifteen Lemmas, Monoid Collapse, Golden Strassen-Clifford 16-MAC sub-cubic tensor execution, and exact algebraic replacements for $\sin, \cos, \tan, \sinh, \cosh, \tanh, \exp,$ and $\log$.
\end{abstract}

\tableofcontents
\newpage

% =============================================================================
\section{System Architecture and Function-to-Theory Master Map}
% =============================================================================

The Zeta Energy System models wave propagation, spinor phase synchronization, and hyperbolic compression across a $3\times 3\times 3$ cubic voxel lattice of spacing $C_{\text{space}} = 25.0$ units (core half-width $C_0 = 12.5$ units). Whereas the legacy 2025 version relied on heuristic floating-point trigonometric interpolation and treated frame $360$ (Step 4) as a spatial singularity, the upgraded 2026i architecture grounds every calculation in the exact algebraic representatives of Hoard (2026i).

Table \ref{tab:master-map} maps every shader routine and JavaScript function in \texttt{zeta\_energy\_system.html} to the exact mathematical structure and theorem in the \textit{Geometric Engine of Time} series that governs it.

\begin{longtable}{@{}p{3.6cm}p{3.8cm}p{4.2cm}p{3.4cm}@{}}
\caption{Master mapping between codebase functions, mathematical equations, and Geometric Engine of Time components.} \label{tab:master-map} \\
\toprule
\textbf{Code Function / Shader} & \textbf{Mathematical Object} & \textbf{Role in Simulation} & \textbf{Engine Source} \\
\midrule
\endfirsthead
\toprule
\textbf{Code Function / Shader} & \textbf{Mathematical Object} & \textbf{Role in Simulation} & \textbf{Engine Source} \\
\midrule
\endhead
\bottomrule
\endfoot
\texttt{getZetaPrimeWheel(t)} \newline \texttt{getZeta(t)} & $Z(t) = \sum_{p \in P_{10}} p^{-1/2}\cos(t\ln p)$ & Prime-weighted harmonic carrier wave across the 10 lowest wheel generators. & Hoard (2026e) Prop.~3.2; Prediction 2 (30-Wheel) \\
\addlinespace
\texttt{vertexShaderTube} & $E_{\text{scalar}} = u_{\text{Scale}} G_{\text{env}}(s_k)\left(\frac{s}{3} + S(t)\right)$ & Modulates tube radial expansion by the $C_3$ apex invariant $s/3$, soliton spikes, and spread $s_k$. & Prop.~2.3(3), Prop.~3.4, Prop.~8.2 \\
\addlinespace
\texttt{fragmentShaderTube} \newline \texttt{fragmentShaderShell} & $l_j = \operatorname{smoothstep}\!\left(\sin(\phi + \frac{2\pi j}{3})\right)$ & Renders the three $120^\circ$ $C_3$ phase lobes and highlights $k \in \{2,6\}$ handshake apices. & Prop.~2.1, Prop.~3.7, Prop.~7.1 \\
\addlinespace
\texttt{GAUSSIAN\_CLOCK[k]} & $\epsilon_k = \frac{(1+i)^k}{2^{\lfloor k/2\rfloor}} \in \{-1,0,1\}^2$ & Exact 8-step projective spinor clock and rational spread $s_k \in \{0, \frac{1}{2}, 1\}$. & Prop.~2.2, Prop.~2.4 (Niven), Prop.~8.2 \\
\addlinespace
\texttt{applyWheelRotation40} & Circulant $R(\rho_0, \rho_1, \rho_2) \in \mathbb{Q}(2\cos 40^\circ)$ & Rotates the base cube triad $\Delta_0$ by $40^\circ$ and $80^\circ$ about $(1,1,1)$ to form $\Delta_1, \Delta_2$. & Prop.~3.1 (Corrected Frame), Prop.~8.3 \\
\addlinespace
\texttt{getCoarseGearBlock} & $A_l = \{m : \lfloor 10m/72\rfloor = l\}$ & Maps fine frame $n \in [0,720]$ to middle gear $m \in [0,71]$ and coarse gear $l \in [0,9]$. & Hoard (2026b); Prop.~3.9 (Pump) \\
\addlinespace
\texttt{ExactTrefoilCurve} & $T(t) = (\sin t + 2\sin 2t, \cos t - 2\cos 2t, -\sin 3t)$ & Embeds the $D_3$-symmetric knot inside the $\{4,3,5\}$ Klein cell ($r_{\max} = 0.2520 < a$). & Prop.~3.11, Section 4.3 \\
\addlinespace
\texttt{calculateNodePos} & $\mathbf{x}_i(n) \in \mathbb{R}^3$ for $i \in \{0..11\}$ & Computes the 12-node trajectory via exact spinor-hyperbolic interpolation or legacy breath. & Prop.~2.3, Prop.~4.3, Prop.~4.4 \\
\addlinespace
\texttt{updateCompressionNodes} & $\mathbf{W}_{A,i}, \mathbf{W}_{B,i} \in \mathbb{R}^3$ & Positions Slice A and Slice B in the $3\times 3\times 3$ grid and updates 3D inter-slice bridges. & Prop.~3.3, Section 4.3 \\
\addlinespace
\texttt{drawSpinorClockDial} & $(k_A, k_B) = (n \bmod 8, n_B \bmod 8)$ & Interactive 8-node dial monitoring dual-slice spinor sheet lock and $\pm i$ handshakes. & Prop.~2.3(2)--(3), Prop.~3.7 \\
\addlinespace
\texttt{drawSlicer} & $d_i^2 = \|\mathbf{x}_{A,i} - \mathbf{x}_{B,i}\|^2$, $\langle X_A, X_B\rangle$ & Detects 2D/3D trajectory hotspots, null-cone inner products, and $3\times 3$ adjacent flux. & Prop.~3.3 (Axiom C2), Prop.~4.3 \\
\addlinespace
\texttt{updateCharts} & $s = a+b+c$, $|z|^2$, $E_{\text{bind}} = 2(ab+bc+ca)$ & Real-time 12-channel spectrum of the $C_3$ Fourier split, quadrances, and cross-terms. & Prop.~3.4, Lemma 5 Part 6(a) \\
\addlinespace
\texttt{updateQuaternionDisplay} & Form 3 64-bit Einstein Tile + Tracks A, B, C, R & Packs branch coordinates into 64 bits and verifies $\bmod\ 881$ residue fingerprints. & Def.~5.3, Prop.~8.7, Prop.~8.8, Def.~8.10 \\
\end{longtable}

\subsection{Integration of the Fifteen Lemmas and Four-Layer Deductive System into Section 1}
By integrating \textit{The Fifteen Lemmas of the Geometric Engine} (Hoard, October 16, 2026), the architectural table above is grounded in a strict **Four-Layer Deductive System plus Layer V (System Level)**:
\begin{enumerate}
    \item \textbf{Origin-Centered Cubit Convention (Global Convention 1--5):} All spatial coordinates in \texttt{zeta\_energy\_system.html} are normalized to the standardized local acoustic cubit
    \begin{equation}
    C_0(f_{\text{cal}}; t, p, x_w, x_c) := \frac{\lambda}{16} = \frac{v_{\text{local}}(t,p,x_w,x_c)}{16 f_{\text{cal}}} = 53.748\text{ cm} \pm 0.022\text{ cm} \quad (\text{at } f_{\text{cal}} = 40\text{ Hz}, 20^\circ\text{C}, 50\%\text{ RH}).
    \end{equation}
    In the simulation, \texttt{ HALF\_CUBE = 12.5} corresponds to $1\,C_0$, so the centered inner core is $K = [-C_0, C_0]^3$ (Cube 1, edge $2C_0 = \lambda/8$, containing the $8$ core voxel centers) and the outer voxel cube is $V = [-2C_0, 2C_0]^3$ (Cube 2, edge $4C_0 = \lambda/4$, containing all $4\times 4\times 4 = 64$ voxel centers). As proven in Global Convention 4 and Lemma 5 Part 2, centering at the origin preserves the full 24-element octahedral rotation group $O \subset SO(3)$ and its 48-element spinor double cover $2O \subset S^3$ (which contains all 24 Hurwitz units $\mathcal{H}_1$, including $h = \frac{1-i+j+k}{2}$), whereas a corner-anchored cube $[0, 4C_0]^3$ preserves only 3 rotations and breaks $h$.
    \item \textbf{Mapping the 5 Toroidal Tiers and the $7.536\,C_0$ Separation Gap (Lemmas 4, 7, 10):} Beyond the Lower Manifold $V = [-2C_0, 2C_0]^3$ (which lies inside radius $\|x\| \le 2\sqrt{3}\,C_0 \approx 3.4641\,C_0$), the Upper Manifold $\mathcal{M}_{\text{upper}} = \bigcup_{n=0}^4 T_n$ consists of five smoothly embedded tori with major radii $R_n = (5n^2 + 15n + 15)C_0 \in \{15, 35, 65, 105, 155\}C_0$, minor radii $r_n = (2n+4)C_0 \in \{4, 6, 8, 10, 12\}C_0$, and integer windings $W_n = 16 - 2n \in \{16, 14, 12, 10, 8\}$, satisfying the exact invariants:
    \begin{equation}
    r_n + W_n = 20\,C_0 = \frac{5}{4}\lambda, \qquad \sum_{n=0}^4 W_n = W_\Sigma = 60, \qquad \operatorname{dist}(V, T_0) \ge (11 - 2\sqrt{3})C_0 \approx 7.5359\,C_0 \approx 4.050\text{ m}.
    \end{equation}
    By Lemma 7 Part 2, the central tier $T_2$ ($r_2 = 8, W_2 = 12$) is the unique arithmetic mean and median of both sequences.
    \item \textbf{Synthetic Trigonometric Elimination at Layer I (Lemmas 1, 2, 3):} Before any shader or quaternion operation is called, Zhao's Three Lemmas eliminate inverse trigonometric calls ($\arccos, \arcsin, \operatorname{atan2}$) from geometric setup:
    \begin{itemize}
        \item \textit{Lemma 1 (Isotomic Tangency Conservation):} Replaces angle-bisector tangent calculations on $\triangle ABC$ with linear semiperimeter differences $BD = CX = s - b$ and $BX = CD = s - c$, where $s = (a+b+c)/2$. This mirrors the mirror-symmetric tier radii $r_n$ and $r_{4-n}$ about $r_2 = 8C_0$.
        \item \textit{Lemma 2 (Spiral Similarity Center $z_0$):} Replaces polar angle/radius similarity solves between paired segments $AB \mapsto CD$ with a single complex rational division:
        \begin{equation}
        \alpha = \frac{c - d}{a - b}, \qquad z_0 = \frac{ad - bc}{a - b - c + d}.
        \end{equation}
        \item \textit{Lemma 3 (Symmedian Quadrance Ratio \& Projective Harmonic Bundle):} Replaces the Law of Sines ($\sin\angle BAD / \sin\angle CAD = \sin C / \sin B$) with the exact ratio of adjacent squared lengths (quadrances $Q_3 = AB^2, Q_2 = AC^2$):
        \begin{equation}
        \frac{BF}{FC} = \frac{EB}{EC} = \frac{Q_3}{Q_2}, \qquad EA^2 = EB\cdot EC, \qquad (B,C; E,F) = -1, \qquad MB^2 = MC^2 = MF\cdot ME.
        \end{equation}
    \end{itemize}
\end{enumerate}

% =============================================================================
\section{Signal Generation: The 30-Wheel Prime Proposer and \texorpdfstring{$C_3$}{C3} Fourier Split}
% =============================================================================

\subsection{The Prime-Weighted Dirichlet/Zeta Wave Function}
Each of the 9 longitudinal streams in the $XY$-grid (spanning $x, y \in \{-25, 0, +25\}$ and $z \in [-30, +30]$) is driven by a truncated Dirichlet-type harmonic series evaluated over the ten fundamental generator primes/units $P_{10}$ associated with the 30-node primorial wheel (Hoard 2026e, Propositions 3.2 and 3.3; Hoard 2026i, Prediction 2):
\begin{equation}
P_{10} = \{1, 2, 3, 5, 7, 11, 13, 17, 19, 23\}.
\end{equation}
In both the GLSL vertex shader (\texttt{getZetaPrimeWheel}) and the JavaScript spectrum engine (\texttt{getZeta}), the real carrier wave at phase coordinate $t$ is computed as:
\begin{equation}
Z(t) = \sum_{p \in P_{10}} \frac{1}{\sqrt{p}} \cos(t \ln p).
\label{eq:zeta-wave}
\end{equation}
For a point at normalized axial coordinate $u_x \in [0, 1]$ along a tube of length $L_{\text{strip}} = 60$ at fine sequencer frame $n \in [0, 720]$, the local phase argument $t$ is:
\begin{equation}
t(u_x, n) = 50 u_x - 8\tau(n), \quad \text{where} \quad \tau(n) = \frac{n}{50}.
\end{equation}

\subsection{Threefold \texorpdfstring{$C_3$}{C3} Branch Offsets and Soliton Spike Injection}
By Proposition 2.1 (The Nine-Phase Factorization), the nine phases of $\mu_9$ decompose into three $120^\circ$ ($2\pi/3$) triads, which are cosets of $\mu_3 = \{1, \omega, \omega^2\}$ with $\omega = e^{2\pi i/3}$. In the shader and spectrum analyzer, the three branches $j \in \{0, 1, 2\}$ of a stream are separated by exact $120^\circ$ carrier shifts scaled by $10\text{ rad}^{-1}$:
\begin{equation}
\Delta t_{120^\circ} = 10 \cdot \frac{2\pi}{3} \approx 20.943951, \qquad \Delta t_{240^\circ} = 10 \cdot \frac{4\pi}{3} \approx 41.887902.
\end{equation}
Thus, at any axial slice $t$, the three real branch amplitudes $(a, b, c)$ scaled by wave amplitude $A = u_{\text{Scale}}$ are:
\begin{equation}
a(t) = A \cdot Z(t), \qquad b(t) = A \cdot Z(t + 20.94395), \qquad c(t) = A \cdot Z(t + 41.88790).
\label{eq:three-branches}
\end{equation}
When the user triggers \texttt{injectBtn} (\textit{Inject Soliton Spike}), a Gaussian pulse timestamp $t_{\text{spike},m} = -\frac{8n}{50} - 5.0$ is pushed into the 10-slot uniform buffer \texttt{uSpikes[10]}, adding a localized soliton excitation:
\begin{equation}
S(t) = \sum_{m=0}^{9} 4 \exp\!\left(-0.2\,(t - t_{\text{spike},m})^2\right).
\end{equation}

\subsection{The \texorpdfstring{$C_3$}{C3} Fourier Split and the 12 Spectrum Analysis Charts}
Proposition 3.4 of Hoard (2026i) proves that the representation $\mathbb{R}^3$ of the cyclic group $C_3$ decomposes into a 1-dimensional fixed apex line and a 2-dimensional rotating plane ($\mathbb{R}^3 \cong \mathbb{R} \oplus \mathbb{R}^2$):
\begin{equation}
s = a + b + c \in \mathbb{R}, \qquad z = a + \omega b + \omega^2 c \in \mathbb{C},
\end{equation}
with exact inverse reconstruction:
\begin{equation}
a = \frac{1}{3}(s + z + \bar{z}), \qquad b = \frac{1}{3}(s + \omega^2 z + \omega \bar{z}), \qquad c = \frac{1}{3}(s + \omega z + \omega^2 \bar{z}).
\end{equation}
Using the cyclotomic identity $1 + \omega + \omega^2 = 0$ (so $\omega + \omega^2 = -1$), the squared magnitude of the rotating plane vector $z$ reduces purely to rational polynomial combinations of $(a,b,c)$ without complex numbers or square roots:
\begin{align}
|z|^2 &= (a + \omega b + \omega^2 c)(a + \omega^2 b + \omega c) \nonumber \\
&= a^2 + b^2 + c^2 + (\omega + \omega^2)(ab + bc + ca) \nonumber \\
&= (a^2 + b^2 + c^2) - (ab + bc + ca).
\label{eq:z-mag-sq}
\end{align}
The function \texttt{updateCharts(time, amp)} evaluates 12 algebraic observables across 41 spatial samples $t \in [-8\tau, 50 - 8\tau]$ to populate Charts \texttt{c1}--\texttt{c12}:
\begin{enumerate}
    \item \textbf{Charts \texttt{c1}, \texttt{c2}, \texttt{c3} (Branch Amplitudes):} $a(t)$, $b(t)$, and $c(t)$ from Eq.~\eqref{eq:three-branches}.
    \item \textbf{Charts \texttt{c4}, \texttt{c5}, \texttt{c6} (Self-Quadrances):} $Q_0 = a^2$, $Q_1 = b^2$, and $Q_2 = c^2$.
    \item \textbf{Charts \texttt{c7}, \texttt{c8}, \texttt{c9} (Pairwise Cross-Interaction Terms):}
    \begin{equation}
    I_{ab} = 2ab, \qquad I_{bc} = 2bc, \qquad I_{ca} = 2ca.
    \end{equation}
    \item \textbf{Chart \texttt{c10} ($C_3$ Apex Scalar Invariant):} $s = a + b + c$ (fixed under cyclic branch permutation $(a,b,c) \mapsto (c,a,b)$).
    \item \textbf{Chart \texttt{c11} (Rotating Plane Norm-Squared):} $|z|^2 = a^2 + b^2 + c^2 - ab - bc - ca$ from Eq.~\eqref{eq:z-mag-sq}.
    \item \textbf{Chart \texttt{c12} (Inter-Lobe Binding Energy):} The interference excess between total squared apex amplitude $s^2 = (a+b+c)^2$ and the sum of individual branch quadrances:
    \begin{equation}
    E_{\text{bind}} = s^2 - (Q_0 + Q_1 + Q_2) = 2(ab + bc + ca) = \frac{2}{3}\left(s^2 - |z|^2\right).
    \end{equation}
\end{enumerate}

\subsection{Integration of Lemmas 4, 5, 11, and 14 \& Trigonometric Replacements for Section 2}
The signal-generation layer of Section 2 connects directly to four results in \textit{The Fifteen Lemmas} and \textit{Golden Strassen-Clifford Toroids}:
\begin{enumerate}
    \item \textbf{Legendre Mod-5 and Garvan Mod-60 Partition of the 30-Node Wheel (Lemma 4 Part 5):}
    The $\phi(30) = 8$ prime-admissible unit residue classes $(\mathbb{Z}/30\mathbb{Z})^\times = \{1, 7, 11, 13, 17, 19, 23, 29\}$ (encoded in Bits $[60:58]$ of the Einstein Tile) split into two 4-element cosets by the Legendre symbol $(b/5)$, which coincides with the level sets of $\chi_{60}^2$ for Garvan's order-4 mock-theta Dirichlet character $\chi_{60}$:
    \begin{equation}
    \mathcal{B}_+ = \ker\!\left(\frac{\cdot}{5}\right) = \{1, 11, 19, 29\} \pmod{30}, \qquad \mathcal{B}_- = \{7, 13, 17, 23\} \pmod{30}.
    \end{equation}
    Splitting the Zeta sum Eq.~\eqref{eq:zeta-wave} into $Z_+(t)$ over $\mathcal{B}_+$ and $Z_-(t)$ over $\mathcal{B}_-$ provides a quadratic-residue parity channel before any primality test is run.
    \item \textbf{Two-Product 128-Bit Lagrange Gram-TQF Replacement for Sine Area and Weitzenb\"ock Ratio (Lemma 4 Part 2 \& Lemma 5 Part 6a):}
    Given the three branch quadrances $Q_1, Q_2, Q_3$ (or 4D difference vectors $u, v \in \mathbb{Z}^4$ with $|u_k|, |v_k| \le 10^6 < 2^{20}$), standard geometry computes triangle area via $\frac{1}{2}\|u\|\|v\|\sin\theta$ or Heron's square-root formula. Lemma 5 Part 6(a) replaces both with Lagrange's overflow-safe 128-bit integer Gram identity:
    \begin{equation}
    A_R = 4\left(\|u\|^2\|v\|^2 - \langle u,v\rangle^2\right) = 4\left(Q_2 Q_3 - \langle u,v\rangle^2\right) = \left(Q_1 + Q_2 + Q_3\right)^2 - 2\left(Q_1^2 + Q_2^2 + Q_3^2\right) < 2^{86}.
    \end{equation}
    By Lemma 4 Part 2, the algebraic identity
    \begin{equation}
    S^2 - 3A_R = 2\left[(Q_1 - Q_2)^2 + (Q_2 - Q_3)^2 + (Q_3 - Q_1)^2\right] = 6\,|z|^2 \ge 0 \quad (S = Q_1 + Q_2 + Q_3)
    \end{equation}
    proves the quadrance form of Weitzenb\"ock's inequality $\chi = \sqrt{A_R}/S \le 1/\sqrt{3}$, reaching its global maximum $\chi = 1/\sqrt{3}$ if and only if the three branch quadrances are balanced ($Q_1 = Q_2 = Q_3$).
    \item \textbf{Golden Soliton Formation vs.\ the Fokker-Planck Circularity Boundary (Lemma 14 \& Golden Strassen-Clifford Paper §4.1):}
    In \textit{Quaternionic Delta Encoding on Golden Strassen-Clifford Toroids}, §4.1 models soliton formation by setting the diffusion tensor $D_2 \to 0$ when colliding streams land on a prime quadrance $Q = N(q) = p$. Lemma 14 Parts 2--4 formalize this mathematically: setting outgoing rates to zero on a target subset $\mathcal{P} \subset \Lambda$ causes probability mass to converge exponentially fast (at rate $\lambda_{\min} > 0$) to the **Harmonic Hitting Measure**:
    \begin{equation}
    H_{x_0}(p) = \left((-L_{TT})^{-1} L_{TP}\right)_{x_0, p},
    \end{equation}
    and on the infinite 4D Hurwitz lattice $\mathcal{H}$, the Gibbs partition function $Z_{\mathcal{H}}(\beta) = 1 + \sum_{n=1}^\infty 24\sigma_{\text{odd}}(n)e^{-\beta V(n)}$ converges whenever $V(n) \ge \alpha \ln n$ with $\beta\alpha > 2$. By Lemma 14 Part 4 (the **Circularity Boundary**), because evaluating the trap indicator $\mathbf{1}_{\mathcal{P}}(Q)$ requires testing whether $Q$ is prime, the Fokker-Planck wave operates in **Helix A** (candidate proposer/sampler) while **Helix B** certifies primality in $O(1)$ using the Form 1 bitmap or Miller--Rabin.
\end{enumerate}

% =============================================================================
\section{The Projective Spinor Clock, Spreads, and the 24-Frame Schedule}
% =============================================================================

\subsection{Why Projective Arithmetic is Mandatory: Niven's Theorem}
A unit quaternion $q = \cos\alpha + u\sin\alpha$ rotates $\mathbb{R}^3$ via $x \mapsto qx\bar{q}$ by the spatial angle $2\alpha$ about axis $u$, where $\alpha$ is the spinor angle (Proposition 2.3). To advance space by $90^\circ$ per step requires a spinor angle of $\alpha = 45^\circ = \pi/4$, which has order $8$.

By Proposition 2.4 (Niven's Bound), a rational unit quaternion of finite order can only have order $1, 2, 3, 4,$ or $6$, because $\cos(2\pi m/n) \in \mathbb{Q}$ occurs solely for $\cos\alpha \in \{0, \pm\frac{1}{2}, \pm 1\}$. Since $\cos 45^\circ = \sqrt{2}/2 \notin \mathbb{Q}$, no rational unit quaternion has order $8$; the spinor step $(1+i)/\sqrt{2}$ lies in the binary octahedral group $2O$ (order 48) rather than the 24 Hurwitz units $2T$.

To eliminate $\sqrt{2}$ and floating-point drift, \texttt{zeta\_energy\_system.html} runs the clock projectively on the Gaussian integer $g = 1+i$ (Proposition 2.2):
\begin{equation}
(1+i)^2 = 2i, \qquad (1+i)^4 = -4, \qquad (1+i)^8 = 16.
\end{equation}
Dividing $(1+i)^k$ by the greatest common divisor of its real and imaginary parts, $\gcd = 2^{\lfloor k/2\rfloor}$, yields the 8 reduced Gaussian states $\epsilon_k = a_k + i b_k \in \{-1, 0, 1\}^2$ stored in \texttt{GAUSSIAN\_CLOCK} (Table \ref{tab:clock-table}).

\begin{table}[htbp]
\centering
\caption{Exact projective Gaussian clock states, spatial rotations, and rational spreads implemented in \texttt{GAUSSIAN\_CLOCK}.}
\label{tab:clock-table}
\begin{tabular}{@{}ccccccc@{}}
\toprule
\textbf{Step $k$} & \textbf{Stage Name} & \textbf{Unreduced $(1+i)^k$} & \textbf{$\gcd = 2^{\lfloor k/2\rfloor}$} & \textbf{Reduced $\epsilon_k$} & \textbf{Spinor $q_k$} & \textbf{Spread $s_k$} \\
\midrule
$0$ & Zero Point & $1$ & $1$ & $1 + 0i$ & $+1$ & $0$ (Node) \\
$1$ & Seed & $1 + i$ & $1$ & $1 + i$ & $(1+i)/\sqrt{2}$ & $1/2$ \\
$2$ & Onload & $2i$ & $2$ & $0 + i$ & $+i$ (Apex) & $1$ (Max) \\
$3$ & Project & $-2 + 2i$ & $2$ & $-1 + i$ & $(-1+i)/\sqrt{2}$ & $1/2$ \\
$4$ & Invert & $-4$ & $4$ & $-1 + 0i$ & $-1$ (Sheet) & $0$ (Node) \\
$5$ & Fold & $-4 - 4i$ & $4$ & $-1 - i$ & $(-1-i)/\sqrt{2}$ & $1/2$ \\
$6$ & Offload & $-8i$ & $8$ & $0 - i$ & $-i$ (Apex) & $1$ (Max) \\
$7$ & Index & $8 - 8i$ & $8$ & $1 - i$ & $(1-i)/\sqrt{2}$ & $1/2$ \\
\bottomrule
\end{tabular}
\end{table}

\subsection{Rational Spreads Without Square Roots}
By Proposition 8.2, the transverse envelope $A_k = |\sin(k\pi/4)|$ is represented exactly by its rational spread:
\begin{equation}
s_k = \frac{b_k^2}{a_k^2 + b_k^2} = \sin^2\!\left(\frac{k\pi}{4}\right) \in \left\{0, \frac{1}{2}, 1\right\}, \qquad c_k = \frac{a_k^2}{a_k^2 + b_k^2} = 1 - s_k.
\end{equation}
In \texttt{vertexShaderTube}, the uniform \texttt{uSpread} passes $s_k$ directly to the GPU to modulate the scalar energy envelope and radial normal displacement:
\begin{equation}
G_{\text{env}}(s_k) = 0.65 + 0.35 s_k, \qquad E_{\text{scalar}} = u_{\text{Scale}} \cdot G_{\text{env}}(s_k)\left(\frac{p_1 + p_2 + p_3}{3} + S(t)\right),
\end{equation}
\begin{equation}
\mathbf{p}_{\text{displaced}} = \mathbf{p} + 0.11\,\mathbf{n}\,E_{\text{scalar}}\left(1 + 0.4\,u_{\text{GoldenBoost}}\right).
\end{equation}
Notice how this corrects the legacy 2025 error (Hoard 2026i, Section 10.2 Errata \#3): Step $k=4$ has spread $s_4 = 0$ (an envelope node on the negative spinor sheet $q_4 = -1$), whereas Steps $k=2$ ($+i$) and $k=6$ ($-i$) are the true handshake apices where $s_2 = s_6 = 1$.

\subsection{The Joint 24-Frame CRT Schedule and \texorpdfstring{$C_3 \times C_2$}{C3 x C2} Handshakes}
In \texttt{updateTelemetryBar(n)}, the sequencer advances the spinor clock and the branch index simultaneously once per frame (Definition 3.6 and Proposition 3.7):
\begin{equation}
k(n) = n \bmod 8, \qquad j(n) = n \bmod 3.
\end{equation}
Because $\gcd(8, 3) = 1$, the Chinese Remainder Theorem yields the group isomorphism:
\begin{equation}
\mathbb{Z}/24\mathbb{Z} \xrightarrow{\sim} \mathbb{Z}/8\mathbb{Z} \times \mathbb{Z}/3\mathbb{Z}, \qquad n \bmod 24 \mapsto (n \bmod 8,\; n \bmod 3).
\end{equation}
Within each 24-frame cycle, the six handshake events ($k \in \{2, 6\}$, where $s_k = 1$) occur at frames $n = 4m + 2$ for $m \in \{0, 1, 2, 3, 4, 5\}$ (Proposition 7.1):
\begin{itemize}
    \item \textbf{Onload Windows ($k = 2$, Spinor $+i$):} Frame $n \equiv 2$ (Branch $j=2$), Frame $n \equiv 10$ (Branch $j=1$), and Frame $n \equiv 18$ (Branch $j=0$).
    \item \textbf{Offload Windows ($k = 6$, Spinor $-i$):} Frame $n \equiv 6$ (Branch $j=0$), Frame $n \equiv 14$ (Branch $j=2$), and Frame $n \equiv 22$ (Branch $j=1$).
\end{itemize}
Every branch $j$ reaches its offload window exactly $\Delta = 12$ frames after its onload window, as $\Delta = 12$ is the unique solution modulo $24$ to $\Delta \equiv 4 \pmod 8$ and $\Delta \equiv 0 \pmod 3$.

\subsection{Integration of Lemmas 4, 5, 11, and 13 \& Trigonometric Replacements for Section 3}
Integrating \textit{The Fifteen Lemmas} enriches the spinor clock and scheduling layer with four exact algebraic mechanisms:
\begin{enumerate}
    \item \textbf{Impossibility of Projective $45^\circ$ Spatial Rotations vs.\ Exact $90^\circ$ Steps (Lemma 13 Part 4):}
    Why does the Engine step space by $90^\circ$ ($\alpha = 45^\circ$ spinor angle) rather than $45^\circ$ ($\alpha = 22.5^\circ$ spinor angle)? Lemma 13 Part 4 proves that for any projective integer quaternion $Q = w + \mathbf{v} \in \mathbb{Z}^4$ rotating $\mathbb{R}^3$ by spatial angle $2\theta$, the ratio $\|\mathbf{v}\|^2 / w^2 = \tan^2\theta$ must be rational. For a $45^\circ$ spatial rotation ($\theta = 22.5^\circ$),
    \begin{equation}
    \tan^2(22.5^\circ) = (\sqrt{2}-1)^2 = 3 - 2\sqrt{2} \notin \mathbb{Q},
    \end{equation}
    proving that a $45^\circ$ spatial step is impossible even projectively over $\mathbb{Z}^4$, whereas $\theta = 45^\circ$ ($90^\circ$ spatial step) has $\tan^2(45^\circ) = 1 \in \mathbb{Q}$, realized by $Q_{\text{step}} = 1+i$.
    \item \textbf{Exact Inter-Tier Signed-Permutation Matrices $U_{\text{even}}, U_{\text{odd}}$ and Hurwitz Spinor $h$ (Lemma 4 Parts 2 \& 4, Prop.~7.12):}
    Across the 5 toroidal tiers $T_0, \dots, T_4$, alternating orientations $\operatorname{Rot}_x(\pi/2)$ and $\operatorname{Rot}_y(\pi/2)$ compose to the exact integer signed-permutation matrices:
    \begin{equation}
    U_{\text{even}} = \begin{pmatrix} 0 & -1 & 0 \\ 0 & 0 & 1 \\ -1 & 0 & 0 \end{pmatrix}, \qquad U_{\text{odd}} = U_{\text{even}}^{-1} = \begin{pmatrix} 0 & 0 & -1 \\ -1 & 0 & 0 \\ 0 & 1 & 0 \end{pmatrix},
    \end{equation}
    lifted to $S^3$ by the order-6 Hurwitz unit $h = \frac{1 - i + j + k}{2}$ ($h^2 = -\bar{h}, h^3 = -1, h^4 = -h, h^5 = \bar{h}, h^6 = +1$). By Proposition 7.12, carrying data out from the core through $T_0 \to T_4$ via $W = U_3 U_2 U_1 U_0$ and offloading back from $T_4 \to T_0$ satisfies $W = I$ and $h\bar{h}h\bar{h} = 1$ with rational radial scale $R_4/R_0 = 155/15 = 31/3$ out and $3/31$ back, and $N(hq) = N(q)$ (zero bit growth).
    \item \textbf{Harmonic Bundle Gate Lengths $(0, 2C_0; 4C_0, \frac{4}{3}C_0) = -1$ (Lemma 5 Part 3):}
    Along the core-to-cube baseline $B = 0, C = 2C_0, E = 4C_0$, the symmedian foot of the Apollonius circle of radius $R_A = 2\sqrt{2}\,C_0$ defines the **Klein-Wave Merge Boundary** $x_F = \frac{4}{3}C_0 = \frac{\lambda}{12}$ (acoustic phase $\pi/6$, quadrance ratio $Q_3/Q_2 = 2$) and radial half-edge conjugate $F_{\text{rad}} = \frac{2}{3}C_0 = \frac{\lambda}{24}$ (phase $\pi/12$), encoded directly in Bits $[49:48]$ of the Einstein Tile Gate Byte.
    \item \textbf{Steinhaus--S\'os Three-Distance Theorem \& Fibonacci Two-Gap Perfection (Lemma 11 Part 3):}
    When the Scale Register (Bits $[35:32]$) selects Fibonacci index $m$ ($N = F_m \in \{5, 8, 13, 21, 34, 55, 89\}$), golden circle subdivision $\{k/\varphi \bmod 1 : 0 \le k < F_m\}$ partitions the circle into $F_{m-1}$ gaps of length $\varphi^{-(m-2)}$ and $F_{m-2}$ gaps of length $\varphi^{-(m-1)}$ (in exact ratio $\varphi : 1$) with zero collisions.
\end{enumerate}

% =============================================================================
\section{The \texorpdfstring{$720 \to 72 \to 10$}{720 -> 72 -> 10} Torsional Markov Pump and Multi-Clock Closures}
% =============================================================================

\subsection{Three-Gear Block Telescoping and Open Holonomy}
The 720-frame master cycle nests 90 spinor cycles ($720 = 90 \times 8$) and 30 joint CRT cycles ($720 = 30 \times 24$). In \texttt{updateTelemetryBar(n)} and \texttt{getCoarseGearBlock(m)}, the frame index $n \in [0, 719]$ drives the three-scale torsional pump of Hoard (2026b) and Proposition 3.9:
\begin{enumerate}
    \item \textbf{Outer 12-Point Fine Gear ($720$ frames):} Tracks fine state increments $\Delta q_n = q_n q_{n-1}^{-1}$ with open boundary convention $q_{-1} = q_0$.
    \item \textbf{Middle 10-Point Aggregate Gear ($72$ blocks):} Indexed by $m = \lfloor n / 10 \rfloor \in \{0, \dots, 71\}$. Each 10-frame block telescopes to its boundary ratio:
    \begin{equation}
    \Delta Q_m = \Delta q_{10m+9}\Delta q_{10m+8}\cdots \Delta q_{10m} = q_{10m+9}q_{10m-1}^{-1}, \qquad \delta_m = \Delta Q_m^{-1}.
    \end{equation}
    \item \textbf{Inner 8-Point Coarse Gear ($10$ blocks):} Indexed by
    \begin{equation}
    l(m) = \left\lfloor \frac{10m}{72} \right\rfloor \in \{0, \dots, 9\},
    \end{equation}
    which partitions the 72 middle blocks into ten contiguous runs $A_l = \{m : \lfloor 10m/72\rfloor = l\}$ of exact sizes:
    \begin{equation}
    (|A_0|, |A_1|, \dots, |A_9|) = (8, 7, 7, 7, 7, 8, 7, 7, 7, 7), \qquad \sum_{l=0}^9 |A_l| = 72.
    \end{equation}
\end{enumerate}
By Proposition 3.9(2), the right-acting middle product $Q^{(10)} = \delta_0 \delta_1 \cdots \delta_{71} = P^{-1}$ and coarse product $Q^{(8)} = \prod_{l=0}^9 \left(\prod_{m \in A_l} \delta_m\right) = P^{-1}$ combine with the open product $P = q_{719}q_0^{-1}$ to yield the pump holonomy:
\begin{equation}
H = Q^{(8)} Q^{(10)} P = P^{-1} P^{-1} P = P^{-1}.
\end{equation}
As proven in Proposition 3.13, if the first increment were defined cyclically ($\Delta q_0 = q_0 q_{719}^{-1}$), the product would telescope to $1$ identically for every sequence of states regardless of numerical drift. The telemetry bar therefore tracks the \textbf{open holonomy} $P = q_n q_0^{-1}$ (\texttt{\#tele-holonomy}), confirming closure ($P = +1$) only when $k = n \bmod 8 = 0$.

\subsection{Physical Frequency Closure Windows}
Under Definition 3.8 ($1\text{ frame} = \frac{1}{60}\text{ s}$), the telemetry pill \texttt{\#tele-closures} evaluates five simultaneous modular divisibility conditions at every frame $n$ (Proposition 3.8 and Table 11):
\begin{align}
\text{Caesium-133 } (9,192,631,770\text{ Hz}): &\quad \frac{\nu_{\text{Cs}}}{60} n = 153,210,529.5\,n \in \mathbb{Z} \iff 2 \mid n, \\
\text{40 Hz Carrier}: &\quad \frac{40}{60} n = \frac{2n}{3} \in \mathbb{Z} \iff 3 \mid n, \\
\text{432 Hz Carrier}: &\quad \frac{432}{60} n = \frac{36n}{5} \in \mathbb{Z} \iff 5 \mid n, \\
\text{Joint 24-Frame Register}: &\quad 24 \mid n, \\
\text{Master Synchronized Closure}: &\quad \operatorname{lcm}(24, 2, 3, 5) = 120 \mid n \quad (6 \text{ closures per } 720 \text{ frames}).
\end{align}

\subsection{Integration of the Torsional Markov Pump Paper, Lemma 7, and Lemma 12 into Section 4}
Combining \textit{The Torsional Markov Memory-Offload Pump} (Hoard, June 21, 2026) with Lemma 7 and Lemma 12 of \textit{The Fifteen Lemmas} yields four critical mathematical upgrades and transcendental replacements:
\begin{enumerate}
    \item \textbf{Why Right-Acting Middle Multiplication is Load-Bearing (Lemma 12 Part 2):}
    In the June 2026 Pump paper (§11), the middle net rotation was written schematically as $Q^{(10)} = \prod_{m=0}^{71} \Delta q_m^{(10)}$. Because quaternion multiplication is non-commutative and inversion reverses product order ($(AB)^{-1} = B^{-1}A^{-1}$), Lemma 12 Part 2 proves that while the outer gear multiplies left-to-right in reverse chronological order $Q^{(12)} = \Delta Q_{71}^{(12\to 10)} \cdots \Delta Q_0^{(12\to 10)}$, the middle and inner gears \textbf{must multiply in right-acting chronological order}:
    \begin{equation}
    Q^{(10)} = \delta_0 \delta_1 \cdots \delta_{71} = \left(\Delta Q_{71}^{(12\to 10)} \cdots \Delta Q_0^{(12\to 10)}\right)^{-1} = \left(Q^{(12)}\right)^{-1}, \qquad Q^{(8)} = \Delta q_0^{(8)} \cdots \Delta q_9^{(8)} = Q^{(10)}.
    \end{equation}
    If $Q^{(10)}$ were multiplied in left-acting order $\delta_{71}\cdots\delta_0$, $Q^{(10)}Q^{(12)} \neq 1$ for generic non-commuting quaternions.
    \item \textbf{Replacing Transcendental SLERP $q^w = \exp(w\log q)$ with Rational Cayley Tangential Lifts (Pump Paper §9 \& §14 vs.\ Lemma 9 Part 3):}
    In §9 of the Pump paper, upward reconstruction from inner to middle ($\Delta_c q_m^{(10)} = (\Delta q_{l(m)}^{(8)})^{w_m}$) and middle to outer ($\Delta_c q_n^{(12)} = (\Delta q_{m(n)}^{(10)})^{1/10}$) uses the transcendental SLERP formula $q^w = \exp(w\log q)$, which §14 quarantines from the executable forward path. To bring upward reconstruction inside the exact rational kernel without $\exp, \log, \arccos,$ or $\sin$:
    \begin{itemize}
        \item Represent any principal-branch unit quaternion $q = (w_0, \mathbf{m}) \in S^3 \cap \mathbb{Q}^4$ ($w_0 > 0$) by its rational Cayley vector $\mathbf{u} = \frac{\mathbf{m}}{1 + w_0} = \tan\!\left(\frac{\theta}{2}\right)\hat{\mathbf{u}} \in \mathbb{Q}^3$ (inverse of $\mathcal{C}_{\text{rat}}(\mathbf{u}) = \frac{(1-\|\mathbf{u}\|^2, 2\mathbf{u})}{1+\|\mathbf{u}\|^2}$, Lemma 9 Part 3).
        \item For small or rational-approximated fractional steps $w \in \mathbb{Q}$ (such as $w = 1/10$), apply the rational tangent-addition / Pad\'e approximant $\mathbf{u}_w = w\,\mathbf{u}\left(1 + \frac{1-w^2}{3}\|\mathbf{u}\|^2\right)^{-1} \in \mathbb{Q}^3$ and map back via $\mathcal{C}_{\text{rat}}(\mathbf{u}_w) \in S^3 \cap \mathbb{Q}^4$. Every lifted increment $\hat{q}_n^{(12)}$ then has exact unit norm $N(\mathcal{C}_{\text{rat}}(\mathbf{u}_w)) = 1$ over $\mathbb{Q}$ with zero floating-point renormalization!
    \end{itemize}
    \item \textbf{Replacing $\epsilon = \log H$ Drift Redistribution with BCH Cross-Products (Pump Paper §11 vs.\ Lemma 7 Part 3b):}
    In §11 of the Pump paper, holonomy drift $\epsilon = \log H$ is redistributed across the three gears via $\Delta q_k^{(r)} \leftarrow \exp(-\lambda_r \epsilon)\Delta q_k^{(r)}$ ($\lambda_{12} + \lambda_{10} + \lambda_8 = 1$). By Lemma 7 Part 3(b), for three drift quaternions $H_a = \exp(\epsilon_a)$ ($\epsilon_a \in \mathfrak{su}(2) \cong \mathbb{R}^3$), the compensating inverse $H_4 = (H_1 H_2 H_3)^{-1} = \bar{H}_3\bar{H}_2\bar{H}_1$ has the Baker--Campbell--Hausdorff expansion:
    \begin{equation}
    -\log H_4 = \epsilon_1 + \epsilon_2 + \epsilon_3 + \epsilon_1 \times \epsilon_2 + \epsilon_1 \times \epsilon_3 + \epsilon_2 \times \epsilon_3 + \mathcal{O}(\|\epsilon\|^3),
    \end{equation}
    allowing non-abelian commutator drift across the three gears to be computed via pure 3D vector cross products and mapped back to $S^3 \cap \mathbb{Q}^4$ via $\mathcal{C}_{\text{rat}}(-\frac{1}{2}\lambda_r \epsilon)$.
    \item \textbf{Exact Rational Valve Kernel for Positive Markov Circulation $J > 0$ (Pump Paper §10 vs.\ Lemma 7 Part 3c \& 3e):}
    Instead of using transcendental exponential rates $e^{-\beta\Delta V}$ to break detailed balance on the gear states $\{12, 10, 8\}$, Lemma 7 Part 3(c) takes a symmetric base walk $P_{ij}^0 = W_{ij}/d_i$ and multiplies each reverse edge $e \in \{(10,12), (8,10), (12,8)\}$ by the exact rational valve factor:
    \begin{equation}
    \rho_e = \frac{1}{1 + \lambda A_e} \in \mathbb{Q} \cap (0, 1], \qquad J = \frac{(1 - \rho_{10,12}\rho_{8,10}\rho_{12,8})\Pi_+}{Z} > 0.
    \end{equation}
    On the 64-voxel torus $\Lambda = (\mathbb{Z}/4\mathbb{Z})^3$ (Lemma 7 Part 3e), translation-invariant rational valves $\rho_k \in (0,1)$ yield the exact uniform stationary distribution $\pi = 1/64$ and positive forward edge current $J_{x \to x+e_k} = \frac{(1-\rho_k)w_k}{64W} > 0$.
\end{enumerate}

% =============================================================================
\section{3D Geometric Substructures and Node Kinematics}
% =============================================================================

\subsection{The \texorpdfstring{$\{4,3,5\}$}{{4,3,5}} Golden Hyperbolic Honeycomb Cell}
Inside the central core cube $K = [-C_0, C_0]^3$ ($C_0 = 12.5$), \texttt{goldenCellGroup} renders the central cell $[-a C_0, a C_0]^3$ of the compact hyperbolic $\{4,3,5\}$ honeycomb in the Klein projective model (Proposition 4.3 and Proposition 8.4):
\begin{equation}
\varphi = \frac{1+\sqrt{5}}{2} \approx 1.61803399, \qquad a = \varphi^{-3/2} = \sqrt{\sqrt{5}-2} \approx 0.48586827.
\end{equation}
Its hyperbolic inradius $r_{\text{in}}$ and face-neighbour Lorentz boost rapidity $\vartheta = 2r_{\text{in}}$ satisfy:
\begin{equation}
r_{\text{in}} = \operatorname{artanh}(a) \approx 0.53063753, \qquad \cosh\vartheta = \varphi, \qquad \sinh\vartheta = \sqrt{\varphi}, \qquad \sinh^2(r_{\text{in}}) = \frac{1}{2\varphi} = \cos 72^\circ.
\end{equation}

\subsection{The Corrected Nine-Phase Tetrahedral Wheel \texorpdfstring{$\mu_9$}{mu9}}
In \texttt{tetraWheelGroup}, three regular tetrahedra $\Delta_0, \Delta_1, \Delta_2$ of side $2\sqrt{2}\,C_0$ share the common apex $\mathbf{x}_* = C_0(1,1,1)$. Following the frame correction of Proposition 3.1(4), the $a=0$ triad lies on the three core cube vertices whose coordinates sum to $-C_0$:
\begin{equation}
\mathbf{v}_{0,0} = C_0(1, -1, -1), \qquad \mathbf{v}_{0,1} = C_0(-1, 1, -1), \qquad \mathbf{v}_{0,2} = C_0(-1, -1, 1).
\end{equation}
To generate the $a=1$ ($+40^\circ$) and $a=2$ ($+80^\circ$) tetrahedra, the function \texttt{applyWheelRotation40(vec, a)} applies the exact circulant rotation matrix $R \in \mathbb{Q}(c)$ where $c = 2\cos 40^\circ$ is a root of $c^3 - 3c + 1 = 0$ (Proposition 8.3):
\begin{equation}
R = \begin{pmatrix} \rho_0 & \rho_1 & \rho_2 \\ \rho_2 & \rho_0 & \rho_1 \\ \rho_1 & \rho_2 & \rho_0 \end{pmatrix}, \qquad \rho_0 = \frac{1+c}{3}, \quad \rho_1 = \frac{3 - c - c^2}{3}, \quad \rho_2 = \frac{c^2 - 1}{3},
\end{equation}
where $\rho_0, \rho_1, \rho_2$ are the three roots of $9y^3 - 9y^2 + 1 = 0$, satisfying $R^T R = I$, $R^3 = g_3$ (the $120^\circ$ coordinate permutation), and $R^9 = I$ exactly.

\subsection{Node Kinematics in \texttt{calculateNodePos(nodeIndex, step)}}
Each slice tracks 12 nodes $i \in \{0, \dots, 11\}$ partitioned into three 4-node branch groups ($i \in \{0..3\}$ Red, $i \in \{4..7\}$ Green, $i \in \{8..11\}$ Blue). Let $\mathbf{e}_i$ denote the 12 midpoints of the cube edges of half-width $C_0 = 12.5$, and let $\mathbf{f}_{f(i)}$ denote the 6 cube face centers where $f(i) = 2\lfloor i/4\rfloor + (i \bmod 2) \in \{0, \dots, 5\}$. The golden Klein face centers are $\mathbf{g}_{f(i)} = a\,\mathbf{f}_{f(i)} = \varphi^{-3/2}\mathbf{f}_{f(i)}$.

In the default \textbf{2026i Exact Spinor-Hyperbolic Mode} (\texttt{exact2026}), the position $\mathbf{x}_i(n)$ at frame $n \in [0, 720]$ is governed by three coupled equations:
\begin{enumerate}
    \item \textbf{Macro-Gear Angle and Transverse Spread Envelope:}
    \begin{equation}
    \Phi(n) = \frac{2\pi n}{720}, \qquad \kappa(n) = \frac{n}{90} \in [0, 8], \qquad s_{\text{env}}(n) = \sin^2\!\left(\frac{\kappa(n)\pi}{4}\right).
    \end{equation}
    Notice that $s_{\text{env}}(n)$ equals $1$ at the macro-handshake frames $n = 180$ ($\kappa = 2$) and $n = 540$ ($\kappa = 6$), and vanishes ($0$) at $n = 0, 360, 720$ ($\kappa = 0, 4, 8$).
    \item \textbf{Hyperbolic Klein Radius Modulation:}
    The radial scale interpolates smoothly between the scaled trefoil knot's maximum Klein radius $r_{\text{trefoil}} = 0.2520$ (Hoard 2026i, Section 4.3), the golden cell boundary $a = 0.4859$, and the unit core cube $1.0$:
    \begin{equation}
    r_{\text{Klein}}(n) = 0.2520 + (1 - 0.2520)\left(0.35 + 0.65\,s_{\text{env}}(n)\right)\left(0.65 + 0.35\cos\Phi(n)\right).
    \end{equation}
    \item \textbf{Edge-to-Golden-Face Lerp and $90^\circ$-per-Octant Spatial Screw:}
    First, the unrotated vector $\mathbf{b}_i(n)$ interpolates between the outer cube edge $\mathbf{e}_i$ and the golden Klein face center $\mathbf{g}_{f(i)}$:
    \begin{equation}
    \mathbf{b}_i(n) = r_{\text{Klein}}(n)\left[\,s_{\text{env}}(n)\,\mathbf{e}_i + \left(1 - s_{\text{env}}(n)\right)\mathbf{g}_{f(i)}\,\right].
    \end{equation}
    Second, the $XY$-plane rotates by $\theta_{\text{sp}}(n, i) = \kappa(n)\frac{\pi}{2} + 0.15\lfloor i/4\rfloor\frac{2\pi}{3}$ (completing $720^\circ$ of spatial rotation across the 8 octants, Proposition 2.3), while the $Z$-coordinate breathes with $\cos(\Phi(n)/2)$:
    \begin{equation}
    \mathbf{x}_i(n) = \begin{pmatrix}
    b_{i,x}(n)\cos\theta_{\text{sp}}(n,i) - b_{i,y}(n)\sin\theta_{\text{sp}}(n,i) \\[0.3em]
    b_{i,x}(n)\sin\theta_{\text{sp}}(n,i) + b_{i,y}(n)\cos\theta_{\text{sp}}(n,i) \\[0.3em]
    b_{i,z}(n)\cos\!\left(\frac{\Phi(n)}{2}\right)
    \end{pmatrix}.
    \label{eq:exact-node-pos}
    \end{equation}
\end{enumerate}

\subsection{Integration of Lemma 8 (Steps 3--6), Ghost Map \texorpdfstring{$G(p)$}{G(p)}, and Rational Trefoil Replacements for Section 5}
Integrating Lemma 8 of \textit{The Fifteen Lemmas} and \textit{Golden Strassen-Clifford Toroids} into Section 5 adds four exact geometric constructions and eliminates all trigonometric calls from the trefoil and hyperbolic projection:
\begin{enumerate}
    \item \textbf{Cube 3 Fibonacci Binomial Cube $W = [-4, 4]^3$ (Lemma 8 Step 5):}
    Continuing the doubling $K = [-1,1]^3 \to V = [-2,2]^3 \to W = [-4,4]^3$ (edge $8C_0 = \lambda/2$, where $8 = F_6 = 5 + 3 = F_5 + F_4$), the Golden Binomial Cube identity $n^3 = (m+r)^3 = m^3 + 3m^2 r + 3m r^2 + r^3$ evaluates in exact integers to:
    \begin{equation}
    8^3 = (5 + 3)^3 = 125 + 225 + 135 + 27 = 512,
    \end{equation}
    partitioning Cube 3 into 1 major cube ($5^3$), 3 slabs ($3 \times 5^2 \times 3$), 3 rods ($3 \times 5 \times 3^2$), and 1 residual cube ($3^3$), while staying separated from $T_0$ by $(11 - 4\sqrt{3})C_0 \approx 4.07\,C_0 \approx 2.19\text{ m}$.
    \item \textbf{The Step 6 Ghost Map $G(p) = \Phi(\mathcal{C}_{\text{rat}}(p))$ and Six Disjoint Ghost Cells (Lemma 8 Propositions 6.2 \& 6.4):}
    For any point $\mathbf{p} \in \mathbb{Q}^3$ inside the open unit ball $|\mathbf{p}| < 1$, the Ghost Map $G = \Phi \circ \mathcal{C}_{\text{rat}}$ maps $\mathbf{p}$ to the upper hyperboloid sheet $\mathbb{H}^3 = \{x_0^2 - |\mathbf{x}|^2 = 1, x_0 > 0\}$ via pure rational operations:
    \begin{equation}
    G(\mathbf{p}) = \left(\frac{1 + |\mathbf{p}|^2}{1 - |\mathbf{p}|^2},\; \frac{2\mathbf{p}}{1 - |\mathbf{p}|^2}\right) \in \mathbb{H}^3 \cap \mathbb{Q}^4.
    \end{equation}
    Applying the six $\{4,3,5\}$ face-neighbour boosts $B_{\pm k}$ ($\cosh\vartheta = \varphi, \sinh\vartheta = \sqrt{\varphi}$, e.g. $B_{+1} = \begin{pmatrix}\varphi & \sqrt{\varphi}\\ \sqrt{\varphi} & \varphi\end{pmatrix} \oplus I_2$) carries the scaled trefoil $G(c T_\varphi)$ ($c = 1/12$) into the six face-neighbouring hyperbolic cubes. By Proposition 7.17 (Collapse at Step 6), at the offload window $k=6$ (spinor $-i$), the inverse boosts $B_{+k}^{-1}$ pull the three positive ghosts back into the central cell where $\Delta_1 \cup (-\Delta_1) = \{\pm 1\}^3$ reconstructs all 8 vertices of the core cube.
    \item \textbf{Rational Weierstrass Substitution $u = \tan(t/2)$ for the $\varphi$-Trefoil (Lemma 8 Prop.~6.1 \& 6.2):}
    In \texttt{ExactTrefoilCurve}, the curve $T_s(t) = (\sin t + 2\sin 2t,\; \cos t - 2\cos 2t,\; -s\sin 3t)$ (with $s = 1$ or $s = \varphi$) normally evaluates six transcendental sines and cosines per point. Setting $u = \tan(t/2) \in \mathbb{Q}$ gives $\cos t = \frac{1-u^2}{1+u^2}$ and $\sin t = \frac{2u}{1+u^2}$, transforming $T_s(t)$ into an **exact rational algebraic curve** over $\mathbb{Q}$ (or $\mathbb{Q}(\sqrt{5})$ for $s = \varphi$):
    \begin{align}
    x(u) &= \frac{2u}{1+u^2} + \frac{8u(1-u^2)}{(1+u^2)^2} = \frac{2u(5 - 2u^2 - 3u^4)}{(1+u^2)^3}, \\
    y(u) &= \frac{1-u^2}{1+u^2} - \frac{2(1 - 6u^2 + u^4)}{(1+u^2)^2} = \frac{-1 + 11u^2 - 5u^4 - u^6}{(1+u^2)^3}, \\
    z(u) &= -s\,\frac{2u(3 - u^2)(1 - 3u^2)}{(1+u^2)^3}.
    \end{align}
    Furthermore, because $T_s(t)$ has $D_3$ dihedral symmetry of order $6$, only the fundamental arc $t \in [0, \pi/3]$ ($u \in [0, 1/\sqrt{3}]$) is needed; the other 5 arcs are generated by the integer signed-permutation matrices $g_3 = \begin{pmatrix}0&1&0\\0&0&1\\1&0&0\end{pmatrix}$ and $g_2 = \begin{pmatrix}0&-1&0\\-1&0&0\\0&0&-1\end{pmatrix} \in O(3,\mathbb{Z})$.
    \item \textbf{Equivalent Golden Clifford Torus Knot Inside $SU(2)$ (Lemma 8 Prop.~6.5 vs.\ Golden Strassen-Clifford Paper §2--3):}
    With the Golden Clifford Toroid radii $r_1 = \frac{\varphi}{\sqrt{\varphi+2}}$ and $r_2 = \frac{1}{\sqrt{\varphi+2}}$ ($r_1/r_2 = \varphi, r_1^2 + r_2^2 = 1$), the curve
    \begin{equation}
    K(s) = \left(r_1 e^{2is},\; r_2 e^{3is}\right) = r_1\cos 2s + r_1\sin 2s\,i + r_2\cos 3s\,j + r_2\sin 3s\,k \in S^3
    \end{equation}
    is the exact $(2,3)$ trefoil knot on the Golden Clifford Torus. As proven in Lemma 11 Part 5, $K(s)$ closes at $s = 2\pi$ because its frequency ratio $2:3$ is rational, whereas choosing an irrational step ratio $\Delta\theta_1 / \Delta\theta_2 = \varphi$ gives the ergodic surface-scanning flow of the Golden Strassen-Clifford paper.
\end{enumerate}

% =============================================================================
\section{Complete Catalog of Interaction Points and Detection Functions}
% =============================================================================

A primary objective of the upgraded dual-slice architecture is to trace two cubic slices---\textbf{Slice A} at voxel $\mathbf{v}_A = (v_{A,x}, v_{A,y}, v_{A,z}) \in \{-1,0,1\}^3$ and depth $Z_A \in [0,1]$, and \textbf{Slice B} at adjacent voxel $\mathbf{v}_B = (v_{B,x}, v_{B,y}, v_{B,z}) \in \{-1,0,1\}^3$ and depth $Z_B \in [0,1]$---and monitor every mathematical point where the two trajectories and their surrounding streams interact.

Given current frame $n \in [0, 720]$ and spinor phase offset slider $\Delta k \in \{0, \dots, 7\}$, Slice B's effective frame index $n_B$ is shifted by $\Delta k$ macro-octants ($90\Delta k$ frames) plus $\Delta k$ fine steps:
\begin{equation}
n_B(n, \Delta k) = (n + 90\Delta k + \Delta k) \bmod 721.
\label{eq:step-b}
\end{equation}
Thus Slice A is in fine spinor state $k_A = n \bmod 8$, and Slice B is in fine spinor state $k_B = (k_A + 2\Delta k) \bmod 8$ (or shifted by $\Delta k$ when stepping via the interactive dial). Below are the \textbf{seven interaction mechanisms}, the exact code functions that compute them, and the Geometric Engine theorems that define them.

\subsection{Interaction Point 1: Spinor Double-Cover Sheet Lock and \texorpdfstring{$\pm i$}{+-i} Handshake Resonance}
\begin{itemize}
    \item \textbf{Code Function:} \texttt{drawSpinorClockDial(step, stepB)} updating \texttt{\#clock-dual-status}.
    \item \textbf{Geometric Engine Component:} Proposition 2.3(2)--(3) and Theorem 7.5 of the \textit{Fifteen Lemmas} (Double-Cover Handshake Apices).
    \item \textbf{Mathematical Condition:}
    Let $\Delta k_{\text{fine}} = (k_B - k_A) \bmod 8$. Three distinct spinor interaction regimes occur:
    \begin{equation}
    \text{Spinor Interaction State} = \begin{cases}
    \text{\textbf{APEX HANDSHAKE }} (\pm i), & \text{if } \{k_A, k_B\} = \{2, 6\} \implies s_{k_A} = s_{k_B} = 1, \; q_{k_B} = -q_{k_A}, \\
    \text{\textbf{SHEET LOCK }} (-q), & \text{if } \Delta k_{\text{fine}} = 4 \implies \operatorname{Rot}_x(90k_A^\circ) = \operatorname{Rot}_x(90k_B^\circ), \; q_{k_B} = -q_{k_A}, \\
    \text{\textbf{PHASE SYNC }} (+q), & \text{if } \Delta k_{\text{fine}} = 0 \implies q_{k_B} = +q_{k_A}.
    \end{cases}
    \end{equation}
    When $\{k_A, k_B\} = \{2, 6\}$, both slices simultaneously achieve maximum transverse spread $s_2 = s_6 = 1$ and identical spatial half-turn geometry $\operatorname{diag}(1, -1, -1)$ on opposite sheets of $SU(2)$, enabling lossless onload-to-offload state transfer across adjacent cubes.
\end{itemize}

\subsection{Interaction Point 2: 3D Inter-Slice Laser Bridge Resonance}
\begin{itemize}
    \item \textbf{Code Function:} \texttt{updateCompressionNodes(step)} updating \texttt{interactionBridgeLines}.
    \item \textbf{Geometric Engine Component:} Section 4.3 ($\{4,3,5\}$ Face-Neighbour Boosts) and Definition 5.3 (Lower Manifold Voxel Lattice $V = [-2C_0, 2C_0]^3$).
    \item \textbf{Mathematical Condition:}
    In 3D world space, the centers of Slice A and Slice B are positioned at:
    \begin{align}
    \mathbf{P}_A &= 25.0\,\left(v_{A,x},\; v_{A,y},\; v_{A,z} + Z_A - 0.5\right), \\
    \mathbf{P}_B &= 25.0\,\left(v_{B,x},\; v_{B,y},\; v_{B,z} + Z_B - 0.5\right).
    \end{align}
    The 12 bridge segments connect world node $\mathbf{W}_{A,i} = \mathbf{P}_A + \mathbf{x}_i(n)$ to $\mathbf{W}_{B,i} = \mathbf{P}_B + \mathbf{x}_i(n_B)$. The 3D bridge interaction detector evaluates the local cube-frame Euclidean distance:
    \begin{equation}
    d_{\text{3D},i}(n, n_B) = \left\|\mathbf{x}_i(n) - \mathbf{x}_i(n_B)\right\|_2.
    \end{equation}
    Whenever $d_{\text{3D},i}(n, n_B) < 6.0\text{ units}$ (i.e., within the golden Klein cell half-width $a C_0 = 0.4859 \times 12.5 = 6.073\text{ units}$!), the bridge vertex colors switch from passive cyan/magenta $(0.13, 0.83, 0.93) \to (0.91, 0.47, 0.98)$ to active golden-emerald resonance $(1.0, 0.85, 0.2) \to (0.1, 0.95, 0.6)$.
\end{itemize}

\subsection{Interaction Point 3: 2D Trajectory Collision Hotspots and Resonance Rings}
\begin{itemize}
    \item \textbf{Code Function:} \texttt{drawSlicer(rotation, amp, time, step)} updating \texttt{\#tele-collisions} and drawing golden collision rings.
    \item \textbf{Geometric Engine Component:} Proposition 3.1 (Tetrahedral Edge Squared Distance $= 8$) and Proposition 4.3.
    \item \textbf{Mathematical Condition:}
    For each of the 12 node pairs $i \in \{0, \dots, 11\}$, \texttt{drawSlicer} computes the squared distance $d_i^2 = \|\mathbf{x}_i(n) - \mathbf{x}_i(n_B)\|_2^2$:
    \begin{enumerate}
        \item \textbf{Hotspot Counter Threshold:} If $d_i^2 < 20.0\text{ units}^2$ ($d_i < 4.4721\text{ units}$), pair $i$ is registered in \texttt{collisionPairs} and increments \texttt{collisionCount} (\texttt{\#tele-collisions}).
        \item \textbf{2D Golden Resonance Zone:} If $d_i < 4.5\text{ units}$ in the superimposed projection, a pulsing golden circle is rendered at the midpoint:
        \begin{equation}
        \mathbf{m}_i = \frac{\mathbf{x}_i(n) + \mathbf{x}_i(n_B)}{2}\cdot \text{scale}, \qquad R_{\text{pulse},i} = \max\!\left(4.0,\; 2.8\,(5.5 - d_i)\right)\text{ px}.
        \end{equation}
    \end{enumerate}
\end{itemize}

\subsection{Interaction Point 4: Conformal Null-Cone Inner Product \texorpdfstring{$\langle X_A, X_B\rangle$}{<XA, XB>}}
\begin{itemize}
    \item \textbf{Code Function:} \texttt{conformalInner} in \texttt{drawSlicer}, displayed in \texttt{\#tele-nullcone}.
    \item \textbf{Geometric Engine Component:} Proposition 3.3 (The Conformal Null Cone in $\mathbb{R}^{4,1}$) and Axiom C2 (Linearization of Metric Constraints).
    \item \textbf{Mathematical Condition:}
    By Proposition 3.3(2), lifting a point $\mathbf{x} \in \mathbb{R}^3$ (normalized by core cube radius $C_0 = 12.5$, so $\hat{\mathbf{x}} = \mathbf{x}/C_0 \in [-1,1]^3$) to the conformal null cone
    \begin{equation}
    X(\hat{\mathbf{x}}) = \left(\hat{\mathbf{x}},\; \frac{1}{2}\|\hat{\mathbf{x}}\|^2,\; 1\right) \in \mathbb{R}^{4,1}
    \end{equation}
    with bilinear form $\langle (\mathbf{x},\alpha,\beta), (\mathbf{y},\alpha',\beta')\rangle = \mathbf{x}\cdot\mathbf{y} - \alpha\beta' - \beta\alpha'$ converts quadratic Euclidean distance into an exact linear inner product:
    \begin{equation}
    \langle X(\hat{\mathbf{x}}_A), X(\hat{\mathbf{x}}_B)\rangle = -\frac{1}{2}\left\|\hat{\mathbf{x}}_A - \hat{\mathbf{x}}_B\right\|_2^2.
    \end{equation}
    In \texttt{drawSlicer}, the ensemble conformal null-cone coupling across all 12 node trajectories is evaluated at every frame as:
    \begin{equation}
    \langle X_A, X_B\rangle_{\text{ensemble}} = -\frac{1}{2 C_0^2}\left(\frac{1}{12}\sum_{i=0}^{11} \left\|\mathbf{x}_i(n) - \mathbf{x}_i(n_B)\right\|_2^2\right).
    \end{equation}
    Because $X(\hat{\mathbf{x}})$ is a null vector ($\langle X(\hat{\mathbf{x}}), X(\hat{\mathbf{x}})\rangle = 0$), $\langle X_A, X_B\rangle_{\text{ensemble}} \le 0$ always, reaching $0$ if and only if the two slices coincide on the null cone.
\end{itemize}

\subsection{Interaction Point 5: Adjacent \texorpdfstring{$3\times 3$}{3x3} Cubic Lattice Energy Flux}
\begin{itemize}
    \item \textbf{Code Function:} \texttt{totalAdjFlux} and the $3\times 3$ mini-heatmap in \texttt{drawSlicer}, displayed in \texttt{\#tele-adjflux}.
    \item \textbf{Geometric Engine Component:} Section 6.1 \& Section 6.3 (Cubic Torus Lattice $(\mathbb{Z}/m\mathbb{Z})^3$ vs.\ Hyperbolic Address Shadow).
    \item \textbf{Mathematical Condition:}
    For Slice A at local phase $t_A = 50 Z_A - 8\tau(n) + 12 v_{A,x} + 18 v_{A,y} + 25 v_{A,z}$, the $3\times 3$ adjacent stream monitor samples all 9 streams $(g_x, g_y) \in \{-1, 0, +1\}^2$:
    \begin{equation}
    E_{g_x, g_y}(t_A) = \left| A \cdot Z(t_A + 14 g_x + 21 g_y) \right|,
    \end{equation}
    and integrates the total adjacent boundary energy surrounding the central $[0,0]$ core stream:
    \begin{equation}
    E_{\text{adj}}(t_A) = \sum_{\substack{g_x, g_y \in \{-1,0,1\} \\ (g_x, g_y) \neq (0,0)}} E_{g_x, g_y}(t_A).
    \end{equation}
\end{itemize}

\subsection{Interaction Point 6: Tangential Bezier Vortex Coupling Between Lobes and Core}
\begin{itemize}
    \item \textbf{Code Function:} Lobe loop in \texttt{drawSlicer} and \texttt{fragmentShaderShell}.
    \item \textbf{Geometric Engine Component:} Proposition 3.4 ($C_3$ Threefold Decomposition) and Hoard (2026a) (Toroidal Orbital Stator Coupling).
    \item \textbf{Mathematical Condition:}
    In the 2D slicer, the three $C_3$ lobes sit at radius $R_{\text{lobe}} = 2.1 \cdot \text{scale}$ at angles $\theta_j = \frac{2\pi j}{3} + \theta_{\text{rot}}(n)$. Each lobe's instantaneous radius swells with the local Zeta energy and active soliton spikes:
    \begin{equation}
    r_j(t) = 0.6\,\text{scale}\,\left(1 + 0.25\,e_j(t)\right), \qquad e_j(t) = A\cdot Z(t + 20.94395 j) + S(t).
    \end{equation}
    When the active CRT branch $j = n \bmod 3$ coincides with a spinor handshake apex ($s_k = 1$, i.e., $k=2$ or $k=6$), lobe $j$ switches to golden apex highlight (\texttt{\#f59e0b}), visually marking the exact frame where energy transfers between the rotating $C_3$ plane $z$ and the central core cube.
\end{itemize}

\subsection{Integration of Layer V Interaction Without Contact (Theorems 7.6--7.10) and Monoid Graph Cost into Section 6}
Section 7.3 of \textit{The Fifteen Lemmas} (\textit{Interaction Without Contact}) and Section 4 of \textit{Monoid Collapse and Mass-Node Traversal} add three non-contact interaction channels between engines that do not share physical points (Theorem 7.6 proves distinct engines in a stack of cell edge $> 334\,C_0$ are strictly disjoint):
\begin{enumerate}
    \item \textbf{Interaction Point 7A — The Five-Moment Vandermonde Channel Across the $4.05\text{ m}$ Gap (Lemma 6 \& Theorem 7.7):}
    Across the radial gap $r \in [2\sqrt{3}\,C_0, 11\,C_0]$ separating $\mathcal{M}_{\text{lower}}$ and $\mathcal{M}_{\text{upper}}$, let $X = r^2/(2q) \in [X_a, X_b]$ and define the five cumulative radial moments of the axial and azimuthal velocity profiles $(U, E)$ with $H = \sqrt{2X}\,E$:
    \begin{equation}
    M(X) = \int_0^X U\,dx, \quad I(X) = \int_0^X H\,dx, \quad J(X) = \int_0^X UH\,dx, \quad S(X) = \int_0^X \left(U^2 - \frac{E^2}{2}\right)dx, \quad C_p(X) = \int_0^X \frac{E^2}{2x}\,dx.
    \end{equation}
    By Lemma 6 Parts 1--4 and Theorem 7.7, the $5\times 5$ generalized Vandermonde matrix $B_{ij} = \int_0^\infty x^{\alpha_i}\beta_j(x)\,dx$ splits into a $2+3$ block Jacobian that is strictly invertible whenever $\alpha \notin \{-1/2, 1/2, 3/2\}$, and Banach's contraction mapping theorem solves $Bc + Q(c,c) = d$ uniquely inside $\|c\| \le 2\|B^{-1}\|\|d\|$. Thus, five scalar numbers $(M, I, J, S, C_p)$ at $X_b$ completely lock the exterior stress $\mathcal{T}_0(X)$ across the gap without contact.
    \item \textbf{Interaction Point 7B — The Projected Hyperbolic Quadrance Channel $Q_h(i,j)$ (Lemma 8 Prop.~6.2 \& Theorem 7.8):}
    When Slice A has spinor state $q_A = \mathcal{C}_{\text{rat}}(\mathbf{p}_A)$ and Slice B has spinor state $q_B = \mathcal{C}_{\text{rat}}(\mathbf{p}_B)$, their non-contact interaction in $\mathbb{H}^3$ is governed by the exact rational kernel:
    \begin{equation}
    Q_h(A, B) = \langle G(\mathbf{p}_A), G(\mathbf{p}_B)\rangle_L^2 - 1 = \sinh^2 d_{\mathbb{H}}\!\left(G(\mathbf{p}_A), G(\mathbf{p}_B)\right),
    \end{equation}
    where $\langle G(\mathbf{p}), G(\mathbf{r})\rangle_L = \frac{(1+|\mathbf{p}|^2)(1+|\mathbf{r}|^2) - 4\mathbf{p}\cdot\mathbf{r}}{(1-|\mathbf{p}|^2)(1-|\mathbf{r}|^2)}$. This replaces transcendental $\operatorname{arccosh}$ and $\sinh$ evaluations with a single rational fraction!
    \item \textbf{Interaction Point 7C — Mass-Node Relational Graph Edge Cost $c_{ij}$ (Monoid Collapse Paper §4.2):}
    In \textit{Monoid Collapse and Mass-Node Traversal}, the relational traversal cost between nodes $i$ and $j$ is $c_{ij} = \frac{d_{ij}^2}{M_i + M_j + \epsilon}$. Integrating Theorem 7.7 and Theorem 7.8 upgrades $d_{ij}^2$ to either the conformal null-cone inner product $-2\langle X(\mathbf{x}_i), X(\mathbf{x}_j)\rangle$ (Euclidean regime) or the hyperbolic quadrance $Q_h(i,j) = \sinh^2 d_{\mathbb{H}}(G_i, G_j)$ (hyperbolic regime), while $M_i, M_j$ are supplied by the first Vandermonde mass moment $M(X_b)$ of Lemma 6:
    \begin{equation}
    c_{ij}^{\text{GET}} = \frac{Q_h(i, j)}{M_i(X_b) + M_j(X_b) + \epsilon} \ge 0.
    \end{equation}
\end{enumerate}

% =============================================================================
\section{64-Bit Einstein Tile Packing and the Four Parallel Tracks}
% =============================================================================

In \texttt{updateQuaternionDisplay(rotation, amp, time, step)}, the live state of the center stream $[0,0,0]$ is converted into a quaternion, packed into a 64-bit Einstein Tile word (Definition 5.3), and cross-checked across the four parallel trigonometric tracks of Section 8.10.

\subsection{Quaternion State Construction Along the \texorpdfstring{$(1,1,1)$}{(1,1,1)} Tetrahedral Axis}
Given the three branch signals $a(t), b(t), c(t)$ and their $C_3$ apex mean $\mu_{\text{apex}} = \frac{a+b+c}{3}$, the deviation quaternion is formed and right-multiplied by the spatial rotation quaternion $q_{\text{rot}}$ about the unit tetrahedral axis $\hat{\mathbf{u}} = \frac{1}{\sqrt{3}}(1,1,1)$:
\begin{equation}
q_{\text{state}} = \mu_{\text{apex}} + (a - \mu_{\text{apex}})i + (b - \mu_{\text{apex}})j + (c - \mu_{\text{apex}})k,
\end{equation}
\begin{equation}
q_{\text{rot}} = \cos\!\left(\frac{\theta_{\text{rot}}}{2}\right) + \frac{i+j+k}{\sqrt{3}}\sin\!\left(\frac{\theta_{\text{rot}}}{2}\right), \qquad q_{\text{final}} = q_{\text{state}} \cdot q_{\text{rot}}.
\end{equation}

\subsection{Bit-Level Encoding into the 64-Bit Einstein Tile (Form 3 Branch Codec)}
Each branch amplitude $v \in \{a, b, c\}$ is quantized into a 10-bit signed two's-complement integer in $[-511, 511]$ via \texttt{pack10Bit(v)}:
\begin{equation}
q_{10}(v) = \operatorname{clamp}\!\left(\operatorname{round}(120\,v),\, -511,\, 511\right) \bmod 1024 \in [0, 1023].
\end{equation}
Let $x_{10} = q_{10}(a)$, $y_{10} = q_{10}(b)$, and $z_{10} = q_{10}(c)$. The 64-bit Einstein Tile is assembled as two 32-bit words (\texttt{upper32} for Bits $[63:32]$ and \texttt{lower32} for Bits $[31:0]$):
\begin{align}
\texttt{lower32} &= \left((j \ \& \ 3) \ll 30\right) \;\mid\; \left((z_{10} \ \& \ \texttt{0x3FF}) \ll 20\right) \;\mid\; \left((y_{10} \ \& \ \texttt{0x3FF}) \ll 10\right) \;\mid\; (x_{10} \ \& \ \texttt{0x3FF}), \\
\texttt{upper32} &= (\texttt{0x1} \ll 30) \;\mid\; \left((j \ \& \ 3) \ll 27\right) \;\mid\; \left((k \ \& \ 7) \ll 24\right) \;\mid\; (\texttt{0xAE} \ll 16) \;\mid\; (\texttt{0x23} \ll 8) \;\mid\; \left(q_{10}(\mu_{\text{apex}}) \ \& \ \texttt{0xFF}\right).
\end{align}
Here:
\begin{itemize}
    \item \textbf{Bits $[29:0]$:} Hold the three 10-bit signed branch coordinates $(z_{10}, y_{10}, x_{10})$.
    \item \textbf{Bits $[31:30]$:} Hold the 2-bit cyclic permutation inverse mode $j = n \bmod 3$ for $B_j^{-1} = g_3^{-j}$ (Open Problem 1).
    \item \textbf{Bits $[39:32]$:} Hold the 8-bit quantized $C_3$ apex invariant $\mu_{\text{apex}}$.
    \item \textbf{Bits $[47:40]$ (\texttt{0x23}):} Nibble \texttt{0x2} selects Track B (exact primary), and nibble \texttt{0x3} specifies Form 3.
    \item \textbf{Bits $[55:48]$ (\texttt{0xAE}):} Gate byte with Holonomy Lock $= 1$ and branch validity mask $[b_2, b_1, b_0] = [1, 1, 1]$.
    \item \textbf{Bits $[63:56]$:} South Halo header encoding the 30-wheel filter class, branch $j \in \{0,1,2\}$, and 3-bit spinor clock step $k \in \{0..7\}$.
\end{itemize}

\subsection{The Four Parallel Verification Tracks (Definition 8.10)}
Finally, the Spectrum panel continuously cross-checks the four parallel trigonometric tracks of Section 8.10:
\begin{enumerate}
    \item \textbf{Track B (Exact Primary, \texttt{\#trk-b}):} Reads the exact reduced Gaussian integer state $\epsilon_k \in \{-1,0,1\}^2$ and rational spread $s_k \in \{0, \frac{1}{2}, 1\}$ of $(1+i)^k$, along with the exact identity $R^9 = I$ in $\mathbb{Q}(2\cos 40^\circ)$.
    \item \textbf{Track A (Floating-Point Readout, \texttt{\#trk-a}):} Displays the spinor angle $\alpha = 45k^\circ$, spatial angle $2\alpha = 90k^\circ$, and $\{4,3,5\}$ hyperbolic cell inradius $r_{\text{in}} = 0.530638$.
    \item \textbf{Track C (Rational Cayley Approximant, \texttt{\#trk-c}):} Applies the rational quaternion $q_C = 138 + 29(i+j+k)$ of norm $N = 138^2 + 3(29^2) = 21,567$ (Proposition 8.7 and Table 13), which approximates the $40^\circ$ wheel rotation to within $0.00116^\circ$ while preserving the circumsphere quadrance $\|\mathbf{x}\|^2 = 3.000000$ exactly over $\mathbb{Q}$.
    \item \textbf{Track R (Triplet Prime Residue Fingerprint, \texttt{\#trk-r}):} Because $p = 881$ satisfies $881 \equiv 1 \pmod 4$, $881 \equiv 1 \pmod 5$, and $881 \equiv -1 \pmod 9$, it splits $\mathbb{Q}(i)$, $\mathbb{Q}(\sqrt{5})$, and $\mathbb{Q}(2\cos 40^\circ)$ simultaneously into 10-bit integer arithmetic ($i \mapsto 387$, $\sqrt{5} \mapsto 228$, $2\cos 40^\circ \mapsto 120, 302, 459 \pmod{881}$, Proposition 8.8). The code computes the modular quadrance fingerprint:
    \begin{equation}
    Q_{881} \equiv \left(x_{10}^2 + y_{10}^2 + z_{10}^2\right) \pmod{881},
    \end{equation}
    certifying the integer payload in $O(1)$ word operations within the Form 2 Garner CRT range $M = 881 \times 883 \times 887 = 690,017,701$.
\end{enumerate}

\subsection{Integration of Section 6.1 (Forms 0--15 Codex Shapes), Lemma 9 Closed-Form Garner Inverses, and Lemma 15 Merkle Security into Section 6}
Integrating Section 6.1, Lemma 9, and Lemma 15 of \textit{The Fifteen Lemmas} completes the bit-level protocol specification:
\begin{enumerate}
    \item \textbf{Exact Closed-Form Polynomial Garner Inverses $c_{12}, c_{23}, c_{13}$ (Lemma 9 Part 2):}
    For any prime triplet $(p, p+2, p+6)$ with $p > 5$ (which must lie in $(11, 13, 17) \pmod{30}$ or $(17, 19, 23) \pmod{30}$ by Lemma 9 Part 1), reconstructing $X \in [0, M)$ from its residues $(x_1, x_2, x_3)$ via
    \begin{equation}
    u_1 = x_1, \quad u_2 = (x_2 - u_1)c_{12} \bmod (p+2), \quad u_3 = ((x_3 - u_1)c_{13} - u_2)c_{23} \bmod (p+6), \quad X = u_1 + u_2 p + u_3 p(p+2)
    \end{equation}
    requires zero extended-GCD calls or lookup tables because all three modular inverses are exact closed-form polynomials in $p$:
    \begin{equation}
    c_{12} = \frac{p+1}{2}, \qquad
    c_{23} = \begin{cases} \frac{3p+17}{4} & p \equiv 1 \pmod 4, \\ \frac{p+5}{4} & p \equiv 3 \pmod 4, \end{cases} \qquad
    c_{13} = \begin{cases} \frac{p+5}{6} & p \equiv 1 \pmod 6, \\ \frac{5p+29}{6} & p \equiv 5 \pmod 6. \end{cases}
    \end{equation}
    Evaluating at $p = 881$ ($881 \equiv 1 \pmod 4$ and $881 \equiv 5 \pmod 6$) yields the exact constants $c_{12} = 441$, $c_{23} = 665$, and $c_{13} = 739$.
    \item \textbf{Complete Form ID Table (Forms 0--15) and Codex Shapes (Section 6.1, Tables 4--6):}
    In the canonical Fifteen Lemmas bit layout (Table 4 of \textit{Fifteen Lemmas}), Bits $[63:61]$ hold the 3-bit cycle step $k$, Bits $[60:58]$ hold the 3-bit wheel unit $u \in \{0..7\}$, Bits $[57:56]$ hold the 2-bit structure selector ($I, A, \bar{A}, E$ or $\mathcal{M}^{(1)}..\mathcal{M}^{(4)}$), Bits $[55:48]$ hold the Gate Byte, Bits $[47:44]$ hold the step-multiplexed index ($s \in \mathbb{Z}/9$, half-cube $l \in \{0..15\}$, ghost index $0..6$, or $\mathfrak{aff}(3)$ coordinate), Bits $[43:40]$ hold the 4-bit Form ID ($0..15$), Bits $[39:36]$ hold the stream slot ($0..15$, forming the 6-bit Stream ID with Bits $[57:56]$), and Bits $[35:32]$ hold the Scale Register ($2^l, F_m,$ or geodesic frequency $\nu$). Forms $4$--$15$ name deterministic geometric codex shapes: Form 4 (64-cell Voxel Cube $V$), Form 5 (8-cell Core $K$), Form 6 (Stella Octangula $\Delta_1 \cup -\Delta_1$), Form 7 (12-point Cuboctahedron), Form 8 (Geodesic Icosahedron of $10\nu^2+2$ points), Form 9 (24 Hurwitz Units), Form 10 (12-element Binary Dihedral Group $2D_3$), Form 11 (9-Ray Circle $x(s)$), Form 12 (7 Hyperbolic $\{4,3,5\}$ Cells), Form 13 ($F_m$ Golden Circle Points), Form 14 (8 Wheel Units mod 30), and Form 15 (Merkle-attested user codex).
    \item \textbf{Recursive Bottom-Up Merkle Attestation vs.\ South Halo Cache Pre-Filter (Lemma 15):}
    Because a 12-layer 64-voxel state has $64^{12} = 2^{72}$ states ($\ge 256$ preimages per 64-bit word) and the 8-bit South Halo has $2^{56}/2^8 = 2^{48}$ payloads per header ID, Lemma 15 proves that the South Halo is strictly an $O(1)$ L1/L2 cache pre-filter, whereas full structural integrity is guaranteed by the bottom-up recursive Merkle chain $\tau_{12}(G) = H(G_{12})$, $\tau_d(G) = H(G_d \parallel \tau_{d+1}(G))$ ($d = 11,\dots,1$) using a 256-bit collision-resistant hash (SHA-256/BLAKE3, $2^{128}$ birthday security).
\end{enumerate}

% =============================================================================
\section{Summary of Interaction Thresholds and Reference Parameters}
% =============================================================================

Table \ref{tab:thresholds} consolidates all numerical constants, geometric radii, and interaction thresholds active in \texttt{zeta\_energy\_system.html}.

\begin{table}[htbp]
\centering
\caption{Summary of exact geometric constants and interaction thresholds in \texttt{zeta\_energy\_system.html}.}
\label{tab:thresholds}
\begin{tabular}{@{}llll@{}}
\toprule
\textbf{Parameter / Threshold} & \textbf{Symbol} & \textbf{Exact / Numerical Value} & \textbf{Governing Proposition} \\
\midrule
Voxel Lattice Spacing & $C_{\text{space}}$ & $25.0\text{ units}$ & Def.~5.3 ($V = [-2C_0, 2C_0]^3$) \\
Core Cube Half-Width & $C_0$ & $12.5\text{ units}$ & Def.~5.3 ($K = [-C_0, C_0]^3$) \\
$\{4,3,5\}$ Klein Cell Half-Width & $a$ & $\varphi^{-3/2} \approx 0.48586827$ & Prop.~4.3 ($a^2 = \sqrt{5}-2$) \\
$\{4,3,5\}$ Cell Physical Half-Width & $a C_0$ & $6.07335\text{ units}$ & Prop.~4.3 \\
$\{4,3,5\}$ Cell Inradius & $r_{\text{in}}$ & $\operatorname{artanh}(a) \approx 0.53063753$ & Prop.~4.3(2), Prop.~8.4(1) \\
Face-Neighbour Boost Rapidity & $\vartheta$ & $2r_{\text{in}} \approx 1.06127506$ & Prop.~4.3(2) ($\cosh\vartheta = \varphi$) \\
Scaled Trefoil Max Klein Radius & $r_{\text{trefoil}}$ & $0.2520 < a$ & Section 4.3, Table 18 \#7 \\
3D Inter-Slice Bridge Threshold & $d_{\text{3D,max}}$ & $6.0\text{ units} \approx a C_0$ & Section 5.2 (Bridge Resonance) \\
3D Collision Hotspot Threshold & $d_{\text{hotspot}}^2$ & $20.0\text{ units}^2$ ($d < 4.472\text{ units}$) & Section 5.3 (\texttt{\#tele-collisions}) \\
2D Golden Resonance Ring Cutoff & $d_{\text{ring}}$ & $4.5\text{ units}$ & Section 5.3 (\texttt{drawSlicer}) \\
Same-Branch Handshake Delay & $\Delta_{\text{CRT}}$ & $12\text{ frames}$ ($\Delta \equiv 4 \bmod 8$) & Prop.~3.7(3) \\
Master Multi-Clock Closure & $N_{\text{sync}}$ & $\operatorname{lcm}(24,2,3,5) = 120\text{ frames}$ & Prop.~3.8, Table 11 \\
Triplet Prime CRT Modulus & $M$ & $881 \times 883 \times 887 = 690,017,701$ & Axiom C5, Prop.~8.8 \\
\bottomrule
\end{tabular}
\end{table}

\subsection{Integration of the Unified Metrological Standardization Table (Section 6.2 of the Fifteen Lemmas) into Section 7}
At the standard calibration frequency $f_{\text{cal}} = 40\text{ Hz}$ in air at $20^\circ\text{C}$, $50\%$ relative humidity, and $101.325\text{ kPa}$ ($v_{\text{local}} = 343.98\text{ m/s}$, wavelength $\lambda = 8.5995\text{ m}$), every simulation threshold in Table \ref{tab:thresholds} translates directly into physical SI dimensions via $1\,C_0 = \lambda/16 = 53.748\text{ cm} \pm 0.022\text{ cm}$ ($dC_0/dT \approx 0.98\text{ mm/K}$):
\begin{itemize}
    \item \textbf{Four-Fold $1/16$ Half-Cube Footprint (Lemma 8 Step 3):} $\frac{1}{16}C_0 = \frac{\lambda}{256} = 3.359\text{ cm}$ (acoustic phase $\pi/128$).
    \item \textbf{Harmonic Offset $MF = x_F - x_M$ (Lemma 5 Part 3):} $\frac{1}{3}C_0 = \frac{\lambda}{48} = 17.916\text{ cm}$ (phase $\pi/24$).
    \item \textbf{Centered Radial Harmonic Conjugate $F_{\text{rad}}$ (Lemma 5 Part 3):} $\frac{2}{3}C_0 = \frac{\lambda}{24} = 35.832\text{ cm}$ (phase $\pi/12$).
    \item \textbf{Klein-Wave Merge Gate $x_F$ (Lemma 5 Part 3):} $\frac{4}{3}C_0 = \frac{\lambda}{12} = 71.664\text{ cm}$ (phase $\pi/6$).
    \item \textbf{Apollonius Circle Radius $R_A = 2\sqrt{2}\,C_0$ (Lemma 5 Part 3b):} $\frac{\sqrt{2}}{8}\lambda = 1.5202\text{ m}$.
    \item \textbf{Cross-Manifold Separation Gap $\operatorname{dist}(V, T_0) = (11 - 2\sqrt{3})C_0$ (Lemma 10 Part 3):} $7.5359\,C_0 = 0.471\lambda = 4.0504\text{ m}$.
    \item \textbf{Outer Engine Enclosure Radius $(R_4 + r_4)C_0 = 167\,C_0$ (Lemma 10 Part 3):} $\frac{167}{16}\lambda = 89.76\text{ m}$.
\end{itemize}

% =============================================================================
\section{Layer V System-Level Integration: The 3:1 FCC Interleave, 12-Engine Memory, Four-Structure Sandwich, and 64 Streams}
% =============================================================================

Section 7 of \textit{The Fifteen Lemmas} derives, without introducing any new lemma, how multiple Geometric Engines stack and operate together as a fault-tolerant, angle-conserving computational fabric. Integrating Layer V into the Zeta Energy System connects the $3\times 3\times 3$ cubic grid to four system-level theorems:

\subsection{The 3:1 Face-Centered Cubic (FCC) Interleave and the Implied Fourth (Theorems 7.1--7.3)}
Let $L = 2\mathbb{Z}^3$ and define the four Klein four-group offsets $o_0 = (0,0,0)$, $o_1 = (1,1,0)$, $o_2 = (1,0,1)$, and $o_3 = (0,1,1)$. By Theorem 7.2:
\begin{enumerate}
    \item The four disjoint translates $L + o_i$ partition the face-centered cubic (FCC) lattice $D_3 = \{x \in \mathbb{Z}^3 : x_1 + x_2 + x_3 \text{ even}\}$.
    \item The $120^\circ$ rotation $g_3$ about $(1,1,1)$ fixes $L + o_0$ and cycles the three face translates $L+o_1 \to L+o_3 \to L+o_2$.
    \item Since $o_1 + o_2 + o_3 \equiv (0,0,0) \pmod 2$, any three translates uniquely determine the **Implied Fourth**—the exact lattice twin of the quaternion group inverse $H_4 = (H_1 H_2 H_3)^{-1}$ (Lemma 7 Part 3b) and the quaternion basis identity $ijk = -1$.
    \item Each point of $D_3$ has $12$ nearest neighbours—the cuboctahedron vectors $(\pm 1, \pm 1, 0)$ and their permutations (Codex Form 7), which split $3 + 6 + 3$ along $(1,1,1)$ and partition by their zero coordinate into three planar groups of four summing to zero (Proposition 7.13).
\end{enumerate}

\subsection{The Twelve-Engine Parity Hierarchy and the 13th Meta-Parity Engine (Propositions 7.11--7.14)}
In the $3\times 3\times 3$ cubic array surrounding the center stream $[0,0,0]$, the 12 edge-midpoint cuboctahedral neighbours form the **Twelve-Engine Hierarchy**:
\begin{equation}
\text{Group } A: \; F_A = (A_1 A_2 A_3)^{-1}, \qquad \text{Group } B: \; F_B = (B_1 B_2 B_3)^{-1}, \qquad \text{Group } C: \; F_C = (C_1 C_2 C_3)^{-1},
\end{equation}
comprising $9$ data engines ($75\%$) and $3$ implied fourth parity engines ($F_A, F_B, F_C$). Adding the central $[0,0,0]$ engine as the **13th Engine** $F_* = (F_A F_B F_C)^{-1}$ (the fourth of the fourths, Proposition 7.11) allows the array to survive the simultaneous loss of any data engine \textit{together with its own group's parity fourth} (e.g., $A_1$ and $F_A$), because $F_*, F_B, F_C$ first reconstruct $F_A = (F_B F_C F_*)^{-1}$, which then reconstructs $A_1 = (A_2 A_3 F_A)^{-1}$! By Proposition 7.14, routing the implied fourths' traffic in the third $120^\circ$ time slot leaves the first two phases purely for data, achieving a $4/3\times$ data-rate gain ($576$ data bits + $192$ parity bits).

\subsection{The Four-Structure Sandwich System \texorpdfstring{$(I, A, \bar{A}, E)$}{(I, A, A-bar, E)} and 64 Logged Streams (Theorems 7.18--7.20)}
Section 7.6 and 7.7 of \textit{The Fifteen Lemmas} organize four 12-engine rings into the **Four-Structure Sandwich** encoded by Bits $[57:56]$ of the Einstein Tile:
\begin{itemize}
    \item $\texttt{00} = I$: Input 12-engine structure ($I_1, \dots, I_{12}$),
    \item $\texttt{01} = A$: Forward Balancer structure ($A_1, \dots, A_{12}$),
    \item $\texttt{10} = \bar{A}$: Inverse Balancer structure ($\bar{A}_j = A_j^{-1}$),
    \item $\texttt{11} = E$: Emergent structure fixed by the slot-wise sandwich closure:
    \begin{equation}
    A_j I_j A_j^{-1} E_j = 1 \iff E_j = A_j I_j^{-1} A_j^{-1}, \qquad j = 1, \dots, 12.
    \end{equation}
\end{itemize}
This sandwich geometry proves three fundamental theorems:
\begin{enumerate}
    \item \textbf{Theorem 7.18 (Balancers Conserve Rotation Angles):} Because conjugation preserves the scalar real part of any quaternion, $\operatorname{Re}(E_j) = \operatorname{Re}(I_j^{-1}) = \operatorname{Re}(I_j)$, so the emergent structure $E$ has the \textit{exact same rotation angles} as the input $I$, about axes rotated by $A_j$.
    \item \textbf{Theorem 7.19 (Nested Cancellation Law):} Forward and inverse quaternion streams cancel identically for all factor values if and only if they are nested palindromically: $q_1 \cdots q_g q_g^{-1} \cdots q_1^{-1} = 1$, whereas non-nested $q_1 \cdots q_g q_1^{-1} \cdots q_g^{-1} \neq 1$.
    \item \textbf{Proposition 7.20 (Automatic Group Closure \& 3-of-8 Erasure Recovery):} Whenever $I_1 I_2 I_3 I_4 = 1$, the emergent group closes automatically in reverse order: $E_4 E_3 E_2 E_1 = A(I_1 I_2 I_3 I_4)^{-1}A^{-1} = 1$. Every loss of up to $3$ of the $8$ engines in a paired group (and $64$ of the $70$ possible 4-engine losses) is recovered bit-exactly!
\end{enumerate}
Together, the 4 structures $\times$ (12 boundary engine streams + 4 interior projection streams $w, x, y, z$) = **$64$ Streams** logging $64 \times 64 = 4,096\text{ bits}$ per cycle step across the $4\times 4\times 4$ voxel cube $V$.

% =============================================================================
\section{Integration of Monoid Collapse and Mass-Node Traversal in Relational Space}
% =============================================================================

The companion report \textit{Monoid Collapse and Mass-Node Traversal in Relational Space: A Hybrid Coarse-Grained Search Framework} formalizes the transition from continuous 3D volumetric diffusion $\partial\rho/\partial t = D\nabla^2\rho$—where the Gaussian point-source kernel $\rho(r,t) = (4\pi Dt)^{-3/2}\exp(-r^2/4Dt)$ spreads as $r(t) \sim \sqrt{Dt}$ and requires evaluating $N_{\text{grid}} \sim (R/h)^3$ grid cells—to a two-stage hybrid search architecture:
\begin{enumerate}
    \item \textbf{Stage 1 (Algebraic Monoid Homomorphism):} A map $\phi: X \to (M, \cdot, e)$ satisfying $\phi(T_1 \circ T_2) = \phi(T_1) \cdot \phi(T_2)$, compressing a $k$-step transformation chain $T = T_1 \circ \cdots \circ T_k$ into an associative product $\phi(T) = \phi(T_1)\cdots\phi(T_k)$.
    \item \textbf{Stage 2 (Sparse Mass-Node Graph Traversal):} A nonnegative weighted graph $G = (V, E, c)$ with gravitational-style edge cost
    \begin{equation}
    c_{ij} = \frac{d_{ij}^2}{M_i + M_j + \epsilon} \ge 0 \qquad (\epsilon > 0),
    \end{equation}
    solved for the optimal path $P^* = \arg\min_{P: s \to t} \sum_{(i,j)\in P} c_{ij}$ via Dijkstra's algorithm in $O((|V|+|E|)\log|V|)$ time whenever $|V| \ll (R/h)^3$ (Theorem 5.1).
\end{enumerate}

\subsection{Resolving the Open Conjectures of the Monoid Collapse Paper Using the Geometric Engine}
In Section 6 (Table 1) of the Monoid Collapse report, two premises were left as \texttt{[HEURISTIC]} or \texttt{[CONJECTURAL]}: (a) constructing an explicit, bounded-cost monoid homomorphism $\phi: X \to M$ that enables fast jump-ahead without bit-width explosion, and (b) grounding the heuristic cost function $c_{ij}$ in exact Engine invariants. \textit{Hoard (2026i)} and \textit{The Fifteen Lemmas} resolve both:

\begin{table}[htbp]
\centering
\caption{Explicit realizations of the Monoid Homomorphism $\phi: X \to (M, \cdot, e)$ supplied by the Geometric Engine.}
\label{tab:monoid-realizations}
\begin{tabular}{@{}p{3.4cm}p{4.0cm}p{3.6cm}p{3.8cm}@{}}
\toprule
\textbf{Monoid $(M, \cdot, e)$} & \textbf{Homomorphism $\phi(T)$} & \textbf{Jump-Ahead Complexity} & \textbf{Governing Lemma / Theorem} \\
\midrule
\textbf{1. Projective Binary Octahedral Group} $M = 2O / \mathbb{R}^\times$ ($|M|=48$) & Maps cube rotations and $(1+i)^k$ spinor steps to $\epsilon_k \in \{-1,0,1\}^2$ and $2T$. & **$O(1)$ word ops** ($g^N = g^{N \bmod n}$, zero bit growth). & Lemma 5 Part 2; Lemma 13 Part 2 \& 4; Thm.~5.1(1) \\
\addlinespace
\textbf{2. Joint CRT Schedule Ring} $M = \mathbb{Z}/8\mathbb{Z} \times \mathbb{Z}/3\mathbb{Z}$ & $\phi(n) = (n \bmod 8,\; n \bmod 3) \in \mathbb{Z}/24\mathbb{Z}$. & **$O(1)$ modular addition**; locks onload/offload at $\Delta = 12$. & Hoard (2026i) Def.~3.6 \& Prop.~3.7 \\
\addlinespace
\textbf{3. Right-Acting Pump Holonomy Monoid} $M = S^3 \cap \mathbb{Q}^4$ & $\phi(T_1 \circ \dots \circ T_{720}) = Q^{(8)}Q^{(10)}Q^{(12)} = q_0 q_{719}^{-1}$. & **$O(1)$ endpoints** for coboundaries; **$O(\log k)$ depth** via prefix scan. & Lemma 12 Parts 1--3; Hoard (2026b) Prop.~2 \\
\addlinespace
\textbf{4. Triplet Prime CRT Ring} $M = \mathbb{Z}/881 \times \mathbb{Z}/883 \times \mathbb{Z}/887$ & $\phi(Q) = (Q \bmod 881, Q \bmod 883, Q \bmod 887)$. & **$O(\log N)$ square-and-multiply** in 30-bit words; $O(1)$ Garner lift. & Lemma 9 Part 2; Hoard (2026i) Prop.~8.8 \\
\bottomrule
\end{tabular}
\end{table}

Furthermore, in the Zeta Energy System's $3\times 3\times 3$ lattice (and the 12-engine hierarchy of Layer V), the mass-node graph $G = (V, E, c)$ is instantiated directly on the $27$ voxel streams (or $12$ cuboctahedral engines), where:
\begin{equation}
d_{ij}^2 = \begin{cases}
-2 C_0^2 \langle X(\hat{\mathbf{x}}_i), X(\hat{\mathbf{x}}_j)\rangle & (\text{Conformal } \mathbb{R}^{4,1} \text{ Euclidean quadrance, Lemma 8 Step 2}), \\[0.4em]
Q_h(i, j) = \langle G(\mathbf{p}_i), G(\mathbf{p}_j)\rangle_L^2 - 1 & (\text{Hyperbolic } \mathbb{H}^3 \text{ Ghost quadrance, Theorem 7.8}),
\end{cases}
\end{equation}
and the node mass $M_i = M_i(X_b) = \int_0^{X_b} U_i(x)\,dx$ is the conserved first Vandermonde moment of Lemma 6 (or the local $C_3$ apex energy $s_i^2 = (a_i + b_i + c_i)^2$).

% =============================================================================
\section{Integration of Golden Strassen-Clifford Toroids and Sub-Cubic 16-MAC Tensor Execution}
% =============================================================================

The fourth paper, \textit{Quaternionic Delta Encoding on Golden Strassen-Clifford Toroids: Wave Collision and Sub-Cubic Matrix Optimization in the SHD-CCP 64-Bit Protocol}, introduces two major components: (1) mapping differential quaternion updates $q_{n+1} = \Delta q_n \cdot q_n$ ($\Delta q = e^{i\Delta\theta_1} e^{j\Delta\theta_2}$) onto the Golden Clifford Torus in $S^3 \subset \mathbb{R}^4$:
\begin{equation}
x^2 + y^2 = r_1^2 = \frac{\varphi^2}{1 + \varphi^2} = \frac{\varphi^2}{\varphi + 2}, \qquad z^2 + w^2 = r_2^2 = \frac{1}{1 + \varphi^2} = \frac{1}{\varphi + 2}, \qquad \frac{r_1}{r_2} = \varphi,
\end{equation}
and (2) exploiting the structured $4\times 4$ matrix representation of quaternion multiplication to reduce naive $4\times 4$ matrix-matrix multiplication from $64$ scalar multiplications ($O(N^3)$) down to **$16$ core multiply-accumulate (MAC) operations** per state step.

\subsection{How the Fifteen Lemmas Rigorize and Optimize the Golden Strassen-Clifford Toroid}
\begin{enumerate}
    \item \textbf{Scalar-Plus-Skew-Symmetric $4\times 4$ Matrix Decomposition (Lemma 8 Step 2.1):}
    Lemma 8 Step 2.1 proves why the $4\times 4$ left-multiplication matrix $Q_a \in M_4(\mathbb{R})$ of a quaternion $q_a = w_a + x_a i + y_a j + z_a k$ requires only $16$ scalar multiplications when multiplying another quaternion state (or $4\times 4$ block):
    \begin{equation}
    Q_a = w_a I_4 + S(x_a, y_a, z_a) = \begin{pmatrix}
    w_a & -x_a & -y_a & -z_a \\
    x_a &  w_a & -z_a &  y_a \\
    y_a &  z_a &  w_a & -x_a \\
    z_a & -y_a &  x_a &  w_a
    \end{pmatrix}, \qquad Q_a^T Q_a = N(q_a)\,I_4, \qquad \det Q_a = N(q_a)^2.
    \end{equation}
    Because $Q_a$ is completely determined by the 4 parameters $(w_a, x_a, y_a, z_a)$ rather than 16 independent matrix entries, multiplying $Q_a \cdot q_b$ executes as a 4-input, 16-MAC directed acyclic graph (DAG). Even better, when $\Delta q_a$ is one of the 8 reduced Gaussian clock states $\epsilon_k \in \{-1, 0, 1\}^2$ (Lemma 13 Part 4) or one of the signed-permutation matrices $U_{\text{even}}, U_{\text{odd}}, g_3, g_2 \in O(3,\mathbb{Z})$ (Lemma 4 Part 2 \& Lemma 8 Step 6), all entries lie in $\{-1, 0, +1\}$, reducing the 16 multiplications to **zero multiplications** (pure bit-sign flips and register swaps)!
    \item \textbf{Replacing Trigonometric Toroidal Coordinates $q(\theta_1, \theta_2)$ with Rational Cayley \& Lucas/Fibonacci Tracks:}
    In §3.1 of the Golden Strassen-Clifford paper, a point on the toroid is parameterized trigonometrically as:
    \begin{equation}
    q(\theta_1, \theta_2) = (r_1\cos\theta_1) + (r_1\sin\theta_1)i + (r_2\cos\theta_2)j + (r_2\sin\theta_2)k.
    \end{equation}
    Using the Exact-First Convention (Hoard 2026i Section 8) and Lemma 8 Proposition 6.5, we replace $\cos\theta_1, \sin\theta_1, \cos\theta_2, \sin\theta_2$ with two rational half-angle parameters $u_1 = \tan(\theta_1/2) \in \mathbb{Q}$ and $u_2 = \tan(\theta_2/2) \in \mathbb{Q}$:
    \begin{equation}
    q(u_1, u_2) = r_1\left(\frac{1-u_1^2}{1+u_1^2} + \frac{2u_1}{1+u_1^2}i\right) + r_2\left(\frac{1-u_2^2}{1+u_2^2}j + \frac{2u_2}{1+u_2^2}k\right),
    \end{equation}
    which satisfies $N(q(u_1, u_2)) = r_1^2 + r_2^2 = \frac{\varphi^2 + 1}{\varphi + 2} = 1$ identically for all $u_1, u_2 \in \mathbb{Q}$ without a single sine, cosine, or floating-point drift!
\end{enumerate}

% =============================================================================
\section{Master Compendium of \texorpdfstring{$\sin, \cos, \tan, \sinh, \cosh, \tanh, \exp, \log$}{sin, cos, tan, sinh, cosh, tanh, exp, log} Replacement Formulas and Optimizations}
% =============================================================================

To make the entire mathematical toolkit directly usable across shaders, hardware pipelines, and analytical proofs, Table \ref{tab:trig-replacements} and Equations \eqref{eq:rep-weierstrass}--\eqref{eq:rep-bch} catalog every transcendental replacement formula and algorithmic optimization established across the four monographs.

\begin{longtable}{@{}p{3.2cm}p{3.8cm}p{4.8cm}p{2.8cm}@{}}
\caption{Master lookup table of transcendental functions replaced by exact algebraic, rational, or integer operations in the Geometric Engine.} \label{tab:trig-replacements} \\
\toprule
\textbf{Legacy Function} & \textbf{Where It Appeared} & \textbf{Exact Geometric Engine Replacement} & \textbf{Source Lemma / Prop.} \\
\midrule
\endfirsthead
\toprule
\textbf{Legacy Function} & \textbf{Where It Appeared} & \textbf{Exact Geometric Engine Replacement} & \textbf{Source Lemma / Prop.} \\
\midrule
\endhead
\bottomrule
\endfoot
$\cos(k\pi/4),$ \newline $\sin(k\pi/4),$ \newline $\tan(k\pi/4)$ & 8-step spinor clock \& transverse envelope $A_k = |\sin(k\pi/4)|$ & Projective Gaussian integer $(1+i)^k / 2^{\lfloor k/2\rfloor} = a_k + i b_k \in \{-1,0,1\}^2$; rational spread $s_k = \frac{b_k^2}{a_k^2+b_k^2} \in \{0, \frac{1}{2}, 1\}$; $\tan^2 = \frac{b_k^2}{a_k^2}$. & Lemma 13 Part 4; Hoard (2026i) Prop.~2.2, 8.2 \\
\addlinespace
$\cos\theta, \sin\theta, \tan\theta$ \newline (general angle) & Unit circle, $S^3$ rotations, Golden Clifford Toroid $q(\theta_1,\theta_2)$ & Rational Cayley / Weierstrass parameter $u = \tan(\theta/2) \in \mathbb{Q}$: $\cos\theta = \frac{1-u^2}{1+u^2}$, $\sin\theta = \frac{2u}{1+u^2}$, $\tan\theta = \frac{2u}{1-u^2}$. & Lemma 9 Part 3; Lemma 8 Prop.~6.2 \\
\addlinespace
$\sin t + 2\sin 2t,$ \newline $\cos t - 2\cos 2t,$ \newline $-s\sin 3t$ & Parametric $(2,3)$-trefoil knot $T_s(t)$ in 3D and $\mathbb{H}^3$ & Rational degree-6 polynomials in $u = \tan(t/2)$ over $[0, \pi/3]$ + 6 integer $D_3$ matrices $g_3, g_2 \in O(3,\mathbb{Z})$. & Lemma 8 Step 6 \& Prop.~6.1--6.2 \\
\addlinespace
$\cos 40^\circ, \sin 40^\circ$ & 9-phase wheel $\mu_9$ \& tetrahedra $\Delta_1, \Delta_2$ & Cubic field $\mathbb{Q}(c)$ ($c^3 - 3c + 1 = 0$) circulant $R(\rho_0, \rho_1, \rho_2)$ ($9y^3 - 9y^2 + 1 = 0$); mod $881$: $c \equiv 120, 302, 459$. & Lemma 8.2; Hoard (2026i) Prop.~8.3, 8.8 \\
\addlinespace
$\frac{\sin\angle BAD}{\sin\angle CAD},$ \newline $\frac{\sin C}{\sin B}$ & Triangle bisectors, symmedians, \& Klein-wave gate $x_F$ & Symmedian \& Apollonius quadrance ratio $\frac{BF}{FC} = \frac{EB}{EC} = \frac{Q_3}{Q_2}$; harmonic bundle $(0,2C_0; 4C_0, \frac{4}{3}C_0) = -1$. & Lemma 3 Parts 1--3; Lemma 5 Parts 1, 3 \\
\addlinespace
$\|u\|\|v\|\sin\theta,$ \newline $\sqrt{s(s-a)(s-b)(s-c)}$ & Triangle area, cross products, \& Weitzenb\"ock $\chi$ & Two-product 128-bit integer Gram quadrea $A_R = 4(Q_2 Q_3 - \langle u,v\rangle^2) < 2^{86}$; $\sin^2\theta = \frac{A_R}{4Q_2 Q_3}$. & Lemma 4 Part 2; Lemma 5 Part 6(a) \\
\addlinespace
$\cos(N\theta), \sin^2(N\theta)$ & Multi-step infinite-order rotations $(3+4i)^N/5^N$ & Chebyshev $T_N(3/5) = a_N/5^N$ and Spread Polynomial recursion $S_N(s) = 2(1-2s)S_{N-1} - S_{N-2} + 2s$. & Lemma 13 Part 1; Hoard (2026i) Prop.~8.6 \\
\addlinespace
$\theta_1 + \theta_2 + \theta_3 = \pi$ & Euclidean flatness vs.\ hyperbolic area defect $\delta$ & Triple Spread Formula $(s_1+s_2+s_3)^2 - 2(s_1^2+s_2^2+s_3^2) - 4s_1 s_2 s_3 = 0$ (Euclidean) or $\sin\delta\prod\sin(\cdots)$. & Hoard (2026i) Prop.~8.9 \\
\addlinespace
$\cosh d_{\mathbb{H}}, \sinh d_{\mathbb{H}},$ \newline $\operatorname{artanh}|\mathbf{p}|$ & Hyperbolic distance \& non-contact kernel in $\mathbb{H}^3$ & Ghost map $G(\mathbf{p}) = \left(\frac{1+|\mathbf{p}|^2}{1-|\mathbf{p}|^2}, \frac{2\mathbf{p}}{1-|\mathbf{p}|^2}\right)$ and rational kernel $Q_h(i,j) = \langle G_i, G_j\rangle_L^2 - 1 = \sinh^2 d_{\mathbb{H}}$. & Lemma 8 Prop.~6.2; Theorem 7.8 \\
\addlinespace
$\cosh\vartheta, \sinh\vartheta,$ \newline $\tanh(\vartheta/2)$ & $\{4,3,5\}$ honeycomb boosts \& Fibonacci torus boosts & Exact golden algebraic ring $\mathbb{Z}[\sqrt{\varphi}]$ ($\cosh\vartheta = \varphi, \sinh\vartheta = \sqrt{\varphi}, \tanh\frac{\vartheta}{2} = \varphi^{-3/2}$) \& Lucas integers $L_n^2 - 5F_n^2 = 4(-1)^n$. & Lemma 8 Prop.~6.4; Hoard (2026i) Prop.~8.4, 8.5 \\
\addlinespace
$q^w = \exp(w\log q),$ \newline $\epsilon = \log H$ & Torsional Pump upward SLERP \& holonomy correction & Rational Cayley tangent scaling $\mathcal{C}_{\text{rat}}(\mathbf{u}_w)$ and BCH cross-product series $-\log H_4 = \sum \epsilon_a + \sum_{a<b}\epsilon_a \times \epsilon_b$. & Lemma 7 Part 3(b); Lemma 9 Part 3 \\
\addlinespace
$e^{-\beta(V(b)-V(a))}$ & Markov pump \& lattice transition probabilities & Rational valve kernel $\rho_e = \frac{1}{1+\lambda A_e}$ and Sparsemax projection on $Q_h$ logits. & Lemma 7 Part 3(c), 3(e); Theorem 7.9 \\
\bottomrule
\end{longtable}

\subsection{Explicit Algebraic Replacement Equations for Direct Implementation}

\paragraph{1. Complete Rational Substitution for $\sin, \cos, \tan$ and 3D Rotations (Lemma 9 Part 3):}
For any rational tangent half-angle vector $\mathbf{u} = (u_x, u_y, u_z) = \tan(\theta/2)\hat{\mathbf{u}} \in \mathbb{Q}^3$ with quadrance $Q_u = \|\mathbf{u}\|^2 = u_x^2 + u_y^2 + u_z^2$:
\begin{equation}
\cos\theta = \frac{1 - Q_u}{1 + Q_u}, \qquad \sin^2\theta = \frac{4 Q_u}{(1 + Q_u)^2}, \qquad \tan^2\theta = \frac{4 Q_u}{(1 - Q_u)^2},
\label{eq:rep-weierstrass}
\end{equation}
and the exact rational $SO(3,\mathbb{Q})$ rotation matrix acting on $\mathbf{x} \in \mathbb{Q}^3$ without square roots is:
\begin{equation}
\operatorname{Rot}(\mathbf{u})\mathbf{x} = \mathbf{x} + \frac{2}{1 + Q_u}\left(\mathbf{u} \times \mathbf{x} + \mathbf{u} \times (\mathbf{u} \times \mathbf{x})\right).
\end{equation}

\paragraph{2. Exact Symmedian and Harmonic Bundle Replacement for Sine Ratios (Lemma 3 \& Lemma 5):}
Given triangle vertex $A$ and baseline $BC = [0, 2C_0]$ with side quadrances $Q_3 = \|A - B\|^2$ and $Q_2 = \|A - C\|^2$, the internal symmedian foot $x_F$ and external Apollonius tangent $x_E$ are computed with zero trigonometry:
\begin{equation}
x_F = \left(\frac{Q_3}{Q_2 + Q_3}\right)2C_0, \qquad x_E = \left(\frac{Q_3}{Q_3 - Q_2}\right)2C_0, \qquad (0, 2C_0; x_E, x_F) = -1.
\end{equation}

\paragraph{3. Exact Hyperbolic Ghost Kernel and BCH Logarithm Replacement (Lemma 7 \& Lemma 8):}
For any two Poincaré ball states $\mathbf{p}, \mathbf{r} \in \mathbb{Q}^3$ ($|\mathbf{p}|, |\mathbf{r}| < 1$) and three small drift vectors $\epsilon_1, \epsilon_2, \epsilon_3 \in \mathbb{Q}^3$:
\begin{equation}
Q_h(\mathbf{p}, \mathbf{r}) = \left(\frac{(1 + |\mathbf{p}|^2)(1 + |\mathbf{r}|^2) - 4\mathbf{p}\cdot\mathbf{r}}{(1 - |\mathbf{p}|^2)(1 - |\mathbf{r}|^2)}\right)^2 - 1, \qquad
\epsilon_{\text{net}} = \epsilon_1 + \epsilon_2 + \epsilon_3 + \epsilon_1\times\epsilon_2 + \epsilon_1\times\epsilon_3 + \epsilon_2\times\epsilon_3.
\label{eq:rep-bch}
\end{equation}

\end{document}